\chapter{绪论}

\section{研究背景与意义}

在量化投资研究中，股票收益的形成机制具有显著的多因素驱动特征。无论是价值、成长、质量、动量、波动率、流动性等经典风格因子，还是宏观环境、行业景气度、公司事件与市场情绪等扩展信息，都会在不同程度上影响股票的定价与未来收益\textsuperscript{[13][14][15]}。随着数据采集技术、量化平台和研究数据库的发展，研究者已经不再局限于少量核心指标，而是可以同时获取财务报表、日频 K 线、技术指标、行业分类、指数成分变化、分红公告、重大事项乃至外部宏观变量等多源信息。由此带来的直接结果是：股票研究对象正在从“少量因子解释收益”逐步走向“高维因子联合刻画收益”，而高维、多源、长时间窗口已经成为当前量化研究的重要特征。

在这样的研究背景之下，股票因子数据整体存在较高的复杂度，这类复杂度主要体现在三个不同的层面。首先是横截面层面，在同一个交易时间节点，研究需要同时统计和监测几十只、数百只甚至更多的股票标的，不同股票之间，还存在行业属性、市值规模、交易风格和价格联动的各类关联特征。其次是因子层面，单只股票在单个交易时刻，会对应多项基本面、技术面以及拓展类的信息因子，各类因子之间的关联程度普遍较高，很容易出现数据冗余和共线性的问题。最后是时间层面，股票因子数据并不是固定不变的，会跟随政策环境、市场流动性、宏观经济周期以及市场情绪的变动，产生阶段性的偏移和变化\textsuperscript{[10][14][17]}。由此能够看出，股票因子研究的基础数据对象，已经不再是简单的二维表格数据，而是自带“股票—因子—时间”三重结构的高维数据体系。如果直接忽略这类立体结构，只将数据当作普通二维矩阵开展分析，就很难完整还原金融市场里真实存在的多维耦合关联\textsuperscript{[6][9][15]}。

传统多因子研究中，更常见的做法是先将原始数据展平为“样本$\times$特征”的矩阵，再使用主成分分析、逐步回归、相关系数筛选、聚类或打分模型进行降维与选股\textsuperscript{[13][14]}。这类方法具有计算简单、工程实现成熟、结果便于解释等优点，因此长期以来一直是实务和学术研究中的主流路径。然而，其隐含前提通常是：样本之间相对独立、特征之间的关系较稳定、时间维度只作为附加标签存在。对于真实股票因子问题而言，这些前提并不总是成立。首先，同一类因子的有效性可能在牛熊转换、货币宽松或风险偏好变化阶段出现明显差异；其次，不同行业和不同样本池中的股票，对相同因子的敏感度并不一致；再次，当因子数量持续扩张时，仅依靠二维矩阵压缩容易把横截面差异、因子共振关系和时间状态变化混合在一起，从而损失原始结构信息，并削弱模型的解释性和稳定性\textsuperscript{[6][7][9][15]}。

高维因子引发的各类问题，既会给模型搭建工作带来阻碍，也会让数据统计和运算处理面临诸多难题；因子的数量不断增多后，多重共线性、噪声不断堆积、模型过拟合、参数波动不稳定等各类问题都会持续加重，所以现阶段的高维因子模型研究，大多重点关注如何在留存核心数据信息的基础上，合理把控数据的潜在维度\textsuperscript{[15][16]}。除此之外，如果沿用传统的单一二维数据框架来处理多维数据，很容易在数据矩阵转换的过程中，破坏数据原本的原始结构；这就会造成降维结果虽然在数值层面可以正常使用，但对应的实际经济意义会大幅弱化。Lettau在高维因子模型的相关研究中提到，经过张量化处理的高维数据，能够在大幅压缩数据维度的同时，保留数据的主要变异特征，提取出的潜在数据成分，也普遍具备清晰的经济解释逻辑\textsuperscript{[15]}。这一点可以证明，在金融研究场景当中，数据降维的核心不只是简单降低数据维度，更关键的是能否在降维完成后，完整保留数据原有结构和对应的研究解释价值。

张量分解方法能够较好化解这类研究矛盾，为建模分析提供了更加贴合数据特征的研究思路；张量可以看作是普通矩阵的高阶拓展形式，能够直接依托数据原本的多模态结构完成特征学习，不用在研究初期就对多维数据进行粗略的扁平化处理\textsuperscript{[4][6][7]}。常见的张量分解方法主要包含CP分解与Tucker分解两类，其中CP分解依靠多个秩一张量的叠加形式还原原始张量数据，比较适合提取数据中通用的潜在特征；Tucker分解依靠核张量和多个因子矩阵相乘的形式，拟合不同模态之间复杂的相互作用，整体的适用范围和灵活度更高。现有的各类研究能够证实，张量方法除了可以应用在图像识别、数据压缩、结构性降维等常规领域，也慢慢被运用到金融时间序列预测、股票波动规律分析、市场预判研究以及股票关联结构识别等金融研究场景中\textsuperscript{[8][9][10][11][12][17]}。在股票因子的相关研究里，张量分解有着较高的应用价值，该方法不仅可以实现高维因子的数据压缩处理，还能同步输出股票、因子、时间三种模态的潜在特征信息，能够为后续识别股票联动规律、挖掘因子之间的共振联系、分析市场时间状态的动态变化，提供统一的研究支撑。

对中国 A 股市场而言，这种方法价值尤为明显。A 股样本本身具有指数成分动态变化、停牌与复牌并存、财务披露存在时滞、市场风格轮动明显、行业结构差异较大等特点。如果仍然依赖静态、二维和单窗口式的分析框架，就很容易把样本池边界变化、时间漂移和多因子共振这些问题割裂开来处理。相比之下，以三阶张量形式组织股票因子数据，能够更自然地把样本边界、因子集合与时间窗口纳入同一表达体系，也更有利于后续进行跨样本池比较与滚动窗口稳健性分析。换言之，张量方法不仅是一个新的“算法工具”，更对应了一种更适合 A 股复杂数据结构的研究组织方式。

从理论研究的角度来讲，本文将张量分解方法运用到股票因子降维和数据模式挖掘的研究当中；该研究可以在金融数据分析的场景中，进一步验证高阶张量建模方式的实际适用效果，也能对比CP分解、Tucker分解和PCA三种研究方法，在结构留存、信息提取以及结果解读方面的优劣与适用范围\textsuperscript{[4][9][15][16]}；同时，本次研究也将原本多用于信号处理、图像分析领域的张量分解方法，更加明确地运用在股票因子研究领域，为解决高维多因子体系，如何在保留原始结构信息的基础上完成高效数据压缩这一问题，提供更为系统的方法参考。当研究不再只聚焦简单的数据预测，还同时兼顾因子的实际解释性、股票数据的结构特征以及不同时段的市场状态时，张量方法自带的多模态数据表征优势，就能够体现出更加突出的理论研究价值。

从实践应用的角度来看，如果可以在数量庞大的高维股票因子数据中，筛选出一小部分信息含量更高、结构逻辑更清楚的低维核心数据特征，就能够有效简化选股模型的整体复杂程度，缓解因子之间的多重共线性问题，也能让因子风格解释和市场风险监控的结果变得更加稳定\textsuperscript{[14][15]}。除此之外，张量分解方法可以输出对应的时间维度特征，依靠这些特征能够识别市场状态的转变、投资风格的轮换以及特殊市场阶段的结构变化；这不仅可以让每一次实验得到的结果更加容易解读，也能帮助研究人员搭建一套稳定性更强、复用性更高的量化研究流程\textsuperscript{[10][11][12][17]}。所以，依托A股市场的真实样本开展股票因子的张量建模、数据降维和特征模式挖掘相关研究，既可以在研究方法上做出一定的创新，也具备十分显著的现实应用意义。

\section{国内外研究现状}

张量分解的基础数学思路，最早来源于Hitchcock针对多元数组表达形式的相关研究；Hitchcock在1927年提出了高阶张量的基础表达形式，也就是将高阶张量拆解为多个乘积项相加的结构，这也为后续规范化的张量分解研究打下了早期的数学基础\textsuperscript{[1]}。在这之后，Tucker在三模因子分析的研究过程中提出了核张量的相关结构，让高阶数据能够借助不同模态的潜在因子矩阵完成对应表达\textsuperscript{[2]}；Harshman则从多模态因子分析的研究方向出发，完善并推广了PARAFAC方法，也让CP分解慢慢成为高阶数据结构分析里的经典研究模型\textsuperscript{[3]}。Kolda和Bader撰写的综述文献，完整梳理并总结了张量分解的数学定义、算法架构以及各类适用场景，是目前张量分解研究领域认可度很高的基础性文献\textsuperscript{[4]}。国内学者胡胜龙针对张量分解唯一性展开的研究，进一步明确了CP分解的可识别条件和理论适用范围，为高维张量模型的结果解释提供了更加充足的理论依据\textsuperscript{[5]}。

从现有研究的方法综述层面来看，夏虹在整理张量方法与实际应用的综述文献中提到，张量建模最主要的优势，就是可以直接处理各类多维结构数据，不用在建模工作刚开始时，就对多模态数据信息进行粗略的扁平化处理\textsuperscript{[6]}。Rabanser、Shchur和Günnemann三人的综述研究，从机器学习的研究视角，系统性介绍了CP分解、Tucker分解两种方法，以及它们在主题模型、关系数据、时序问题等场景中的具体用途；该研究也提出，张量分解已经不再是单纯的数学分析工具，已经逐步发展为高维数据分析中适用性较广的一类通用方法\textsuperscript{[7]}。综合这些现有研究能够发现，张量方法的核心优势并不只是体现在处理高维数据上，更大的价值在于它可以完整保留数据的原有结构。

在金融数据分析方向，国内已有研究开始尝试将张量模型引入具体金融任务，而不再满足于把高维数据简单展平成二维矩阵。岳欣的研究以停牌情形下的股票波动分析为切入点，选取 2018 年 6 月 1 日至 2020 年 12 月 31 日上证 50 成分股中多个停牌股票的多项交易指标，将其组织为张量形式，并采用基于 CP 分解的 CP-WOPT 算法对缺失值进行联合填充\textsuperscript{[8]}。该研究解决的是“停牌导致样本缺失后如何保留多指标联合结构”的问题，说明张量分解不仅适用于降维，也可用于缺失恢复与后续波动分析；但其任务重心仍是停牌补全和波动刻画，并未进一步讨论降维后的因子排序或结构解释。曾亚丽则把研究对象放在证券市场指数预测上，先利用单个证券市场多元时间序列构造二阶张量，再利用多个证券市场相关技术指标构造三阶张量，并分别引入$(2D)^2$PCA+PCA、MPCA 和 ENW-MPCA 等特征提取算法\textsuperscript{[10]}。其中，$(2D)^2$PCA+PCA 试图解决二维张量特征提取后维数仍然偏高的问题，MPCA 用于解决向量化方法无法保持三维内部结构的问题，而 ENW-MPCA 则进一步针对不同特征向量对预测贡献程度不同这一现象，通过特征值归一化加权提升预测精度。该研究较为系统地说明了张量表示在时间序列预测中的优势，但它服务的是市场指数预测任务，而不是股票因子降维与横截面解释。

从更为通用的统计学习角度进行分析，马少沛等人的相关研究，重点探究了高阶张量如何在不破坏原有结构的前提下完成数据降维的问题\textsuperscript{[9]}；该研究在Zhong等人提出的张量充分降维理念基础上，把原本的二阶研究思路拓展为通用性更强的高阶研究思路，同时以三阶张量作为研究案例，设计出了对应的结构转换算法和结构保持算法，主要用来解决高维张量降维工作中常见的各类问题，包括协方差矩阵奇异、整体计算成本过高、向量化操作损毁原始数据结构等难点。研究人员还将这套新算法和K-means、t-SNE、多维主成分分析、多维判别分析、张量SIR等多种主流方法做了对比测试，结果能够看出，结构保持算法在数据分类的精准度上有着更好的表现。虽然这项研究是以图像识别作为应用场景，但也给本文的研究带来了重要的参考启发，高维张量分析的核心目的不只是简单降低数据维度，更多的是在降维的整个过程中，留存数据原本的结构关联特征。除此之外，Brandi等人尝试将张量分解方法运用到股票结构的相关研究中，提出股票之间的联动规律，不能只依靠二维收益相关系数来衡量；Spelta结合张量分解与链接预测方法，分析了金融市场的可预测特征；Wang等人则以高维向量自回归时序模型为研究基础，证明张量分解可以用于复杂时序系统的参数压缩和潜在结构挖掘工作\textsuperscript{[11][12][17]}。综合以上各类研究可以发现，张量方法在金融领域的应用价值，已经不再局限于单一的数学建模层面，还可以运用在数据修复、结构识别、市场预测、参数压缩等多个研究方向中。

在多因子选股领域，现有中文研究更多聚焦于因子筛选、赋权和宏观修正，而不是显式利用张量结构。黄宏运等从股票选取中的多因子问题出发，收集了 2937 只股票的 10 项影响因子，既包括市净率、市盈率、资产负债比率等基本面指标，也包括 ADR、RSI、BIAS、OBV、RSV 等技术面指标，并以相对收益率为被解释变量建立多元线性回归模型\textsuperscript{[13]}。该研究不仅做了显著性检验，还专门讨论了多重共线性和异方差问题，并进一步采用广义逐步回归法对模型进行改进，使改进后模型的拟合效果明显提升。其贡献在于为多因子选股提供了一套较规范的计量建模流程，但由于方法建立在二维回归框架上，因此难以同时处理股票、因子和时间三个维度的结构耦合。刘宇轩则进一步把宏观金融状态纳入多因子模型，在传统主成分分析法构建多因子选股模型的基础上，引入由 GARCH 模型合成的中国金融状况指数（FCI），构建基于金融周期的行业轮动多因子选股模型\textsuperscript{[14]}。该文不仅比较了金融周期指标与经济周期指标对选股模型的修正效果，还讨论了因子正负性会随市场状态变化而改变，并采用动态回归法判断当期因子正负性。其实证结果表明：引入宏观因素后，多因子模型收益表现得到改善，而金融周期指标相较于经济周期指标更能捕捉金融市场波动，从而取得更高的超额收益率。这说明传统多因子模型的一个关键不足，正是缺少对时间阶段与宏观状态的显式建模。

若把上述中文研究放在同一条问题链中观察，可以发现它们其实分别回答了几个彼此相关、但尚未被统一起来的问题。黄宏运和刘宇轩解决的是“给定因子集合后，如何利用统计回归、主成分分析和宏观修正提升选股效果”的问题\textsuperscript{[13][14]}；岳欣解决的是“当金融数据存在停牌缺失时，如何利用张量分解恢复多指标数据并开展后续分析”的问题\textsuperscript{[8]}；曾亚丽解决的是“当证券市场数据具有多市场、多指标和时间耦合特征时，如何通过张量特征提取提高预测精度”的问题\textsuperscript{[10]}；马少沛等解决的是“高阶张量降维时如何避免结构信息丢失并控制计算复杂度”的问题\textsuperscript{[9]}；李博则通过时态支持向量机和粒度变换，在股票操纵模式发现任务中说明了时间粒度变化本身就可能承载原始特征中难以直接发现的模式信息，其在未知数据集上的识别准确率可达到较高水平\textsuperscript{[18]}。这些研究分别触及了多因子选股、张量建模、时间结构识别和高维降维中的关键环节，但尚未将正式数据链、股票因子压缩、模式发现、排序有效性和样本边界稳健性评价统一进同一套研究框架。换言之，已有研究往往只覆盖“数据处理”“预测建模”或“因子选股”中的某一段链条，而本文希望补上的，正是这几段链条之间缺失的衔接关系。

综合目前已有的相关研究能够发现，张量分解方法在处理高维结构化数据方面，已经被证实具备相应的理论优势，也能够在金融时序数据、股票波动特征以及关联结构建模等研究场景中发挥作用，拥有不错的应用前景；但现有研究依旧存在几处可以进一步完善和探究的地方。首先，以“股票因子—多维结构保序降维—特征模式挖掘”为整体主线的完整研究数量不多，针对国内A股市场样本开展的系统性实证分析也相对匮乏。其次，多数相关研究的重点集中在提升模型预测准确度上，很少同时对因子解释效果、股票联动特征以及市场时间状态切换这三方面内容展开综合分析。最后，很多研究的数据处理流程、实验得出的结果和论文的论证内容之间，缺少统一且清晰的对应逻辑关系\textsuperscript{[8][10][11][18]}。本文的研究便是针对以上这些现存研究缺陷展开，尝试将张量分解研究方法、A股因子数据以及实证结果分析内容整合起来，构建一套完整统一的研究分析框架。

\section{研究内容与技术路线}

本文以“股票—因子—时间”的三维张量结构作为整体研究对象，把整体研究内容划分成四个具体的研究层面。首先是数据层面，本文搭建了一套可重复使用的A股标准化数据处理流程，对K线数据、财务数据、研究报告数据以及指数历史成分数据等基础资料进行整理汇总，最终整理形成完整的因子面板数据。其次是方法层面，本文选用CP分解、Tucker分解和PCA三种不同的降维方式进行对比实验，重点分析张量类研究方法，在数据结构保留和潜在特征模式提取方面的实际优势。再次是结果层面，本文从分解效果质量、排序有效程度、时间状态变动规律以及因子重要程度等多个角度，开展对应的实证分析工作。最后是证据层面，研究将各类数值运算结果、可视化图形结果以及文本论证内容整合在同一个评价体系当中，以此保障整篇论文的研究结论拥有统一、可靠的证据支撑。

对应的技术路线可以概括为“原始数据获取与整理—因子面板构建—三阶张量表示—张量分解与秩选择—潜在表示分析—实验评估与可视化输出”。其中，原始数据整理并不是附属步骤，而是后续张量建模能否成立的前提。由于股票市场数据天然存在停牌、成分股变化、财务披露时滞和不同数据源更新口径不一致等问题，本文在正式实验之前，先对 K 线前复权口径、财务/报告覆盖、指数成分历史和因子面板时间窗口进行统一，确保模型输入具有一致的时间跨度和样本边界。

在算法路线方面，本文先构造原始张量，再执行协议化预处理和训练/验证/测试切分，并分别在 CP、Tucker 和 PCA 三个模型上进行秩选择和留出评估。与仅比较重构误差的实验不同，本文同时关注三个问题：其一，哪种方法在重构质量上更优；其二，哪种方法在排序指标和滚动稳定性上表现更平衡；其三，哪种方法更有利于形成可解释的因子和时间模式。围绕这三个问题，正文给出表格、图形和定性讨论。

\section{论文结构安排}

剔除摘要、参考文献、附录等常规附属内容后，整篇论文的主体内容一共划分成六个章节。第一章为绪论部分，主要介绍本次课题的研究背景、国内外的相关研究现状，同时梳理本文的核心研究内容与整体技术研究路线。第二章为理论基础部分，主要讲解张量代数、张量分解的相关知识，以及金融因子数据表征的基础理论内容。第三章用于说明实验数据的相关信息，包含数据来源、样本选取范围、数据预处理步骤以及三维张量的具体构建方法。第四章详细介绍本文的核心研究方法，也就是依托张量分解实现股票因子降维与数据模式挖掘的整套研究方案。第五章结合完整的实验数据结果，从分解质量、模式识别效果、因子排序有效性等多个维度开展实证分析工作。第六章对全文的研究内容进行整体总结，归纳本次研究得出的结论、研究存在的创新之处、现存的研究缺陷，同时对后续的相关研究方向做出展望。

\chapter{理论基础与关键技术}

\section{张量代数基础}

本文对股票因子数据进行处理时，将其整理为三阶张量$$\mathcal{X}\in\mathbb{R}^{N\times F\times T}$$的形式；其中参数$N$代表研究样本中的股票总数，参数$F$代表选取的因子总数量，参数$T$对应的是整体研究的观测时间节点数量。对比普通的向量和矩阵数据形式，张量最突出的特点，就是可以同时保存多个模态的数据信息，所以该形式更加适配高维结构化数据的描述工作。如果直接把完整的张量数据展平转换为普通矩阵，原本分属不同模态的各类结构特征，就会被强行混合在一起；这会让数据对应的经济含义变得模糊，不利于后续的研究解读。综合以上问题，本文选择直接依托原始的张量数据形态开展建模分析工作。

为便于后文描述 CP 分解与 Tucker 分解，现对本文频繁使用的几类基础运算分别作出说明\textsuperscript{[4][6]}。与上一版只给出压缩公式不同，这里进一步按照“定义—直观解释—示意图—与后文关系”的顺序展开，使各运算之间的逻辑关系更加清楚。

\subsection{模态展开}

模态展开也被称作张量矩阵化，是张量数据处理的基础操作方式；该操作的具体规则，就是固定其中一个模态作为矩阵的行维度，再将剩余的其他模态按照既定顺序展开，作为矩阵的列维度。对于三阶张量$\mathcal{X}\in\mathbb{R}^{N\times F\times T}$来说，如果将股票模态设为矩阵的行，就能得到对应的第一模展开矩阵$X_{(1)}\in\mathbb{R}^{N\times (FT)}$；如果将因子模态设为矩阵的行，就能得到第二模展开矩阵$X_{(2)}\in\mathbb{R}^{F\times (NT)}$；如果将时间模态设为矩阵的行，就可以得到第三模展开矩阵$X_{(3)}\in\mathbb{R}^{T\times (NF)}$。这种操作的本质，是从不同的观测角度对三维数组进行改写，将其转化为普通矩阵，能够把原本复杂的高阶优化问题，转变为矩阵层面的常规计算问题。

在直观上，模态展开可以理解为“沿某一维度观察张量”。当研究股票之间的相对差异时，更关注$X_{(1)}$；当研究因子之间的共振关系时，更关注$X_{(2)}$；当研究时间状态变化时，则更容易从$X_{(3)}$中提取信息。后文 CP 分解的交替最小二乘更新就是在不同模态展开矩阵上轮流进行的，因此模态展开是后续分解算法的入口。

\begin{figure}[H]
\centering
\begin{tikzpicture}[font=\small, >=Latex]
  \matrix (slice1) [matrix of nodes,
    nodes={draw, minimum width=0.8cm, minimum height=0.52cm, anchor=center},
    row sep=-\pgflinewidth, column sep=-\pgflinewidth] {
    $x_{111}$ & $x_{121}$ & $x_{131}$ \\
    $x_{211}$ & $x_{221}$ & $x_{231}$ \\
    $x_{311}$ & $x_{321}$ & $x_{331}$ \\
  };
  \node[above=0.16cm of slice1.north] {时间切片 $t=1$};
  \matrix (slice2) [matrix of nodes,
    nodes={draw, minimum width=0.8cm, minimum height=0.52cm, anchor=center},
    row sep=-\pgflinewidth, column sep=-\pgflinewidth] at ($(slice1.east)+(3.0,0)$) {
    $x_{112}$ & $x_{122}$ & $x_{132}$ \\
    $x_{212}$ & $x_{222}$ & $x_{232}$ \\
    $x_{312}$ & $x_{322}$ & $x_{332}$ \\
  };
  \node[above=0.16cm of slice2.north] {时间切片 $t=2$};
  \node (tensorlabel) at ($(slice1.south)!0.5!(slice2.south)+(0,-0.48)$) {原始三阶张量 $\mathcal{X}$ 的两个切片};

  \matrix (m1) [matrix of nodes,
    nodes={draw, minimum width=0.76cm, minimum height=0.52cm, anchor=center},
    row sep=-\pgflinewidth, column sep=-\pgflinewidth] at ($(tensorlabel.south)+(0,-1.5)$) {
    $x_{111}$ & $x_{121}$ & $x_{131}$ & $x_{112}$ & $x_{122}$ & $x_{132}$ \\
    $x_{211}$ & $x_{221}$ & $x_{231}$ & $x_{212}$ & $x_{222}$ & $x_{232}$ \\
    $x_{311}$ & $x_{321}$ & $x_{331}$ & $x_{312}$ & $x_{322}$ & $x_{332}$ \\
  };
  \node[below=0.2cm of m1.south] {$X_{(1)}$：以股票模态为行};
  \draw[->, thick] ($(tensorlabel.south)+(0,-0.15)$) -- ($(m1.north)+(0,0.22)$);
\end{tikzpicture}
\caption{三阶张量的第一模展开示意图}
\label{fig:mode1-unfolding}
\end{figure}

图\ref{fig:mode1-unfolding}展示出第一模展开的完整实现流程；第一模展开的操作方式，不会修改张量内部各类元素的数值大小，只是把不同时间切片下的因子列数据，按顺序整合排列到同一个矩阵当中，矩阵的每一行数据，依旧对应单一的个股样本。完成这一步处理后，股票模态的相关信息会完整留存于矩阵的行维度当中，因子模态和时间模态的信息则统一整合至矩阵的列维度；后续的算法运算，就能够依托矩阵数据形式，完成股票维度潜在载荷的估算工作。

和第一模展开的处理方式大体相同，第二模展开与第三模展开，仅仅只是调换了用来充当矩阵行的模态类型。在将因子模态设置为行维度时，矩阵里的每一行数据，都代表着单个因子在不同股票以及不同时间节点下的全部数值，这种形式更加便于研究各个因子之间存在的同步变动规律。在将时间模态设置为行维度时，矩阵中的每一行都对应着一个实际观测时间点，能够更好地用来探究市场整体结构随时间产生的各类变化。由此能够看出，这三种不同的张量展开形式，实际上代表着三种不一样的研究分析角度，并不是几种相互独立、没有关联的数据转换方式。

\begin{figure}[H]
\centering
\begin{tikzpicture}[font=\small]
  \matrix (m2) [matrix of nodes,
    nodes={draw, minimum width=0.85cm, minimum height=0.55cm, anchor=center},
    row sep=-\pgflinewidth, column sep=-\pgflinewidth] {
    $x_{111}$ & $x_{211}$ & $x_{311}$ & $x_{112}$ & $x_{212}$ & $x_{312}$ \\
    $x_{121}$ & $x_{221}$ & $x_{321}$ & $x_{122}$ & $x_{222}$ & $x_{322}$ \\
    $x_{131}$ & $x_{231}$ & $x_{331}$ & $x_{132}$ & $x_{232}$ & $x_{332}$ \\
  };
  \node[below=0.18cm of m2.south] {$X_{(2)}$：以因子模态为行};
  \matrix (m3) [matrix of nodes,
    nodes={draw, minimum width=0.75cm, minimum height=0.55cm, anchor=center},
    row sep=-\pgflinewidth, column sep=-\pgflinewidth] at ($(m2.south)+(0,-2.1)$) {
    $x_{111}$ & $x_{211}$ & $x_{311}$ & $x_{121}$ & $x_{221}$ & $x_{321}$ & $x_{131}$ & $x_{231}$ & $x_{331}$ \\
    $x_{112}$ & $x_{212}$ & $x_{312}$ & $x_{122}$ & $x_{222}$ & $x_{322}$ & $x_{132}$ & $x_{232}$ & $x_{332}$ \\
  };
  \node[below=0.18cm of m3.south] {$X_{(3)}$：以时间模态为行};
\end{tikzpicture}
\caption{三阶张量的第二模与第三模展开示意图}
\label{fig:mode23-unfolding}
\end{figure}

\subsection{外积}

若向量$a\in\mathbb{R}^{N}$、$b\in\mathbb{R}^{F}$、$c\in\mathbb{R}^{T}$的外积记为$a\circ b\circ c$，则其生成一个秩一三阶张量，并满足
\begin{equation}
\left(a\circ b\circ c\right)_{ijt}=a_i b_j c_t.
\end{equation}
外积的关键作用在于，它给出了“秩一张量”的最基本构造方式。也就是说，只要给定三个不同模态上的向量，就可以构造出一个在三个模态上同时具有结构的最简单张量。

在后文 CP 分解中，原始张量正是由若干个秩一张量相加而成。因此，理解外积实际上就是理解 CP 分解中“一个潜在成分到底长什么样”。若某一成分的股票向量强调金融股、因子向量强调价值和质量、时间向量强调宽松阶段，那么三者的外积就对应了一个具有明确金融解释的三模态结构。

\begin{figure}[htbp]
\centering
\begingroup
\small
\setlength{\arraycolsep}{18pt}
\renewcommand{\arraystretch}{1.12}
\[
\begin{array}{ccc}
\begin{array}{|c|}\hline a_1\\ \hline a_2\\ \hline a_3\\ \hline\end{array} &
\begin{array}{|c|}\hline b_1\\ \hline b_2\\ \hline b_3\\ \hline\end{array} &
\begin{array}{|c|}\hline c_1\\ \hline c_2\\ \hline\end{array} \\
a & b & c
\end{array}
\]
\vspace{-0.6em}
\[
\left(a\circ b\circ c\right)_{ijt}=a_i b_j c_t
\]
\vspace{-0.4em}
\setlength{\arraycolsep}{10pt}
\[
\begin{array}{cc}
t=1\text{ 切片} & t=2\text{ 切片}\\[0.3em]
\begin{array}{|c|c|c|}
\hline
a_1b_1c_1 & a_1b_2c_1 & a_1b_3c_1\\ \hline
a_2b_1c_1 & a_2b_2c_1 & a_2b_3c_1\\ \hline
a_3b_1c_1 & a_3b_2c_1 & a_3b_3c_1\\ \hline
\end{array}
&
\begin{array}{|c|c|c|}
\hline
a_1b_1c_2 & a_1b_2c_2 & a_1b_3c_2\\ \hline
a_2b_1c_2 & a_2b_2c_2 & a_2b_3c_2\\ \hline
a_3b_1c_2 & a_3b_2c_2 & a_3b_3c_2\\ \hline
\end{array}
\end{array}
\]
\vspace{-0.5em}
\[
a\circ b\circ c\text{ 生成的秩一三阶张量}
\]
\endgroup
\caption{三个向量外积生成秩一三阶张量的示意图}
\label{fig:outer-product}
\end{figure}

\subsection{Kronecker 积}

对于两个矩阵$A=(a_{ij})\in\mathbb{R}^{I\times J}$与$B\in\mathbb{R}^{K\times L}$，其 Kronecker 积记为$A\otimes B$，定义为
\begin{equation}
A\otimes B=
\begin{bmatrix}
a_{11}B & a_{12}B & \cdots & a_{1J}B\\
a_{21}B & a_{22}B & \cdots & a_{2J}B\\
\vdots & \vdots & \ddots & \vdots\\
a_{I1}B & a_{I2}B & \cdots & a_{IJ}B
\end{bmatrix}\in\mathbb{R}^{(IK)\times(JL)}.
\end{equation}
也就是说，Kronecker 积会把矩阵$A$中的每一个元素都替换成“该元素乘以整个矩阵$B$”所形成的分块。它的结果不是按元素相乘，而是把一个小矩阵扩展成一个更大的分块矩阵，因此在结构上更接近“放大并复制模板”的过程。

Kronecker 积本身在本文中不是主角，但它是理解 Khatri--Rao 积的重要前置概念。后者实际上就是“对列向量逐列做 Kronecker 积”，因此若 Kronecker 积不清楚，后面的 Khatri--Rao 积很难真正看懂。

\begin{figure}[htbp]
\centering
\begin{tikzpicture}[font=\small]
  \matrix (A) [matrix of nodes,nodes={draw, minimum width=0.8cm, minimum height=0.6cm},row sep=-\pgflinewidth,column sep=-\pgflinewidth] {
    $a_{11}$ & $a_{12}$ \\
    $a_{21}$ & $a_{22}$ \\
  };
  \node[below=0.2cm of A.south] {$A$};
  \matrix (B) [matrix of nodes,nodes={draw, minimum width=0.8cm, minimum height=0.6cm},row sep=-\pgflinewidth,column sep=-\pgflinewidth] at ($(A.east)+(2.4,0)$) {
    $b_{11}$ & $b_{12}$ \\
    $b_{21}$ & $b_{22}$ \\
  };
  \node[below=0.2cm of B.south] {$B$};
  \draw[->, thick] ($(B.east)+(0.5,0)$) -- ++(1.2,0);
  \matrix (K) [matrix of nodes,nodes={draw, minimum width=0.95cm, minimum height=0.6cm},row sep=-\pgflinewidth,column sep=-\pgflinewidth] at ($(B.east)+(5.0,0)$) {
    $a_{11}b_{11}$ & $a_{11}b_{12}$ & $a_{12}b_{11}$ & $a_{12}b_{12}$ \\
    $a_{11}b_{21}$ & $a_{11}b_{22}$ & $a_{12}b_{21}$ & $a_{12}b_{22}$ \\
    $a_{21}b_{11}$ & $a_{21}b_{12}$ & $a_{22}b_{11}$ & $a_{22}b_{12}$ \\
    $a_{21}b_{21}$ & $a_{21}b_{22}$ & $a_{22}b_{21}$ & $a_{22}b_{22}$ \\
  };
  \node[below=0.2cm of K.south] {$A\otimes B$};
\end{tikzpicture}
\caption{Kronecker 积的分块扩展示意图}
\label{fig:kronecker}
\end{figure}

\subsection{Khatri--Rao 积}

若$A=[a_1,\ldots,a_R]\in\mathbb{R}^{I\times R}$、$B=[b_1,\ldots,b_R]\in\mathbb{R}^{J\times R}$，则 Khatri--Rao 积记为
\begin{equation}
A\odot B=\left[a_1\otimes b_1,\;a_2\otimes b_2,\;\ldots,\;a_R\otimes b_R\right]\in\mathbb{R}^{(IJ)\times R}.
\end{equation}
从定义可以看出，Khatri--Rao 积不是把整个矩阵$A$与整个矩阵$B$做 Kronecker 积，而是把二者的对应列逐列配对，再对每一对列向量做 Kronecker 积，最后把结果重新拼成一个新矩阵。正因为它保留了“列对应”的结构，所以在 CP 分解中非常常见。

对于本文的研究工作来说，Khatri--Rao 积起到的最关键作用，就是将因子模态和时间模态当中的数据信息整合在一起，合并成为一个整体矩阵，方便这类矩阵和各类模态展开矩阵共同完成最小二乘迭代更新的相关计算。简单来说，该运算也是在 CP 分解过程里，用来整合不同模态之间关联信息的主要手段。

\begin{figure}[htbp]
\centering
\begin{tikzpicture}[font=\small, >=Latex]
  \matrix (A) [matrix of nodes,nodes={draw, minimum width=0.8cm, minimum height=0.55cm},row sep=-\pgflinewidth,column sep=-\pgflinewidth] {
    $a_{11}$ & $a_{12}$ \\
    $a_{21}$ & $a_{22}$ \\
  };
  \node[below=0.2cm of A.south] {$A=[a_1,a_2]$};
  \matrix (B) [matrix of nodes,nodes={draw, minimum width=0.8cm, minimum height=0.55cm},row sep=-\pgflinewidth,column sep=-\pgflinewidth] at ($(A.east)+(2.6,0)$) {
    $b_{11}$ & $b_{12}$ \\
    $b_{21}$ & $b_{22}$ \\
    $b_{31}$ & $b_{32}$ \\
  };
  \node[below=0.2cm of B.south] {$B=[b_1,b_2]$};
  \draw[->, thick] ($(B.east)+(0.6,0)$) -- ++(1.2,0);
  \matrix (KR) [matrix of nodes,nodes={draw, minimum width=1.1cm, minimum height=0.55cm},row sep=-\pgflinewidth,column sep=-\pgflinewidth] at ($(B.east)+(5.0,0)$) {
    $a_{11}b_{11}$ & $a_{12}b_{12}$ \\
    $a_{11}b_{21}$ & $a_{12}b_{22}$ \\
    $a_{11}b_{31}$ & $a_{12}b_{32}$ \\
    $a_{21}b_{11}$ & $a_{22}b_{12}$ \\
    $a_{21}b_{21}$ & $a_{22}b_{22}$ \\
    $a_{21}b_{31}$ & $a_{22}b_{32}$ \\
  };
  \node[below=0.2cm of KR.south] {$A\odot B=[a_1\otimes b_1,\;a_2\otimes b_2]$};
\end{tikzpicture}
\caption{Khatri--Rao 积的逐列配对示意图}
\label{fig:khatri-rao}
\end{figure}

\subsection{Hadamard 积}

对于两个同型矩阵$M=(m_{ij}),N=(n_{ij})\in\mathbb{R}^{I\times J}$，其 Hadamard 积记为$M\ast N$，定义为
\begin{equation}
\left(M\ast N\right)_{ij}=m_{ij}n_{ij}.
\end{equation}
Hadamard 积是最直接的按元素乘积。它不会改变矩阵的形状，也不会重新组合行列结构，而是单纯把同一位置上的两个元素对应相乘。因此，相比 Kronecker 积和 Khatri--Rao 积，Hadamard 积在结构上最简单。

在本文后续使用的CP迭代更新公式当中，Hadamard积会应用在格拉姆矩阵的组合部分里，主要用来整合不同模态因子矩阵相互之间的内积相关信息。这种运算的表现形式比较简单，不过在整体求解的过程当中，它能够起到汇总多个模态二阶信息的实际作用。

\begin{figure}[htbp]
\centering
\begin{tikzpicture}[font=\small, >=Latex]
  \matrix (M) [matrix of nodes,nodes={draw, minimum width=0.9cm, minimum height=0.6cm},row sep=-\pgflinewidth,column sep=-\pgflinewidth] {
    $m_{11}$ & $m_{12}$ & $m_{13}$ \\
    $m_{21}$ & $m_{22}$ & $m_{23}$ \\
  };
  \node[below=0.2cm of M.south] {$M$};
  \matrix (N) [matrix of nodes,nodes={draw, minimum width=0.9cm, minimum height=0.6cm},row sep=-\pgflinewidth,column sep=-\pgflinewidth] at ($(M.east)+(2.4,0)$) {
    $n_{11}$ & $n_{12}$ & $n_{13}$ \\
    $n_{21}$ & $n_{22}$ & $n_{23}$ \\
  };
  \node[below=0.2cm of N.south] {$N$};
  \draw[->, thick] ($(N.east)+(0.6,0)$) -- ++(1.2,0);
  \matrix (H) [matrix of nodes,nodes={draw, minimum width=1.15cm, minimum height=0.6cm},row sep=-\pgflinewidth,column sep=-\pgflinewidth] at ($(N.east)+(5.0,0)$) {
    $m_{11}n_{11}$ & $m_{12}n_{12}$ & $m_{13}n_{13}$ \\
    $m_{21}n_{21}$ & $m_{22}n_{22}$ & $m_{23}n_{23}$ \\
  };
  \node[below=0.2cm of H.south] {$M\ast N$};
\end{tikzpicture}
\caption{Hadamard 积的按元素相乘示意图}
\label{fig:hadamard}
\end{figure}

\subsection{与后文分解模型的关系}

上文提到的各类数学运算，并非相互独立的公式组合，它们会在下文构建的张量分解模型中协同发挥作用。模态展开能够将三阶张量转换为矩阵形式，为交替最小二乘算法搭建基础运算条件；外积确定了秩一张量的构建方法，也是构成CP分解最基础的组成部分；Kronecker积是理解列向量耦合拓展关系的基础；Khatri--Rao积能够将两种模态下的列信息逐一进行组合；Hadamard积则用来整合不同模态格拉姆矩阵中逐元素对应的结构信息。把这些运算的作用相互结合进行梳理，就能够更轻松理解后文更新公式里同时出现$X_{(1)}$、$\odot$与$\ast$等符号的原因。

在股票因子的实际应用研究当中，张量代数带来的价值，既能够体现在数据数学表达层面，也可以落实到金融市场的实际含义解读当中。其中股票模态可以用来反映市场横截面的整体结构，因子模态代表各类投资信息对应的维度，时间模态则能够体现出市场整体行情的发展变化趋势。倘若某一项提取得到的潜在成分，在股票维度的载荷数值大多集中在银行、券商以及保险类个股当中，同时在时间维度上又大多出现于政策较为宽松的市场阶段，那么这一潜在成分就可以看作是金融类风格因子在对应市场环境下形成的整体联动走势。由此不难看出，张量划分出的每一类模态都具备清晰明确的经济内涵，并不只是单纯用来划分数据的普通维度。

\section{张量分解模型}

张量分解的核心思路，就是在最大程度留存原始数据多模态结构的基础上，依靠数量更少的潜在成分来近似表达高维张量数据。放到本文研究的股票因子张量当中，该过程既能够实现因子数据的精简压缩，还可以完成个股关联结构以及市场时间状态特征的提取工作。结合现有的经典理论成果以及本文设定的研究目标，后文研究将主要采用CP分解和Tucker分解这两种主流张量分解模型\textsuperscript{[1][2][3][4]}。

\subsection{CP 分解}

CP 分解将原始张量表示为若干个秩一张量之和，其数学形式为
\begin{equation}
\mathcal{X}\approx \sum_{r=1}^{R}\lambda_r a_r\circ b_r\circ c_r.
\end{equation}
其中，$R$为分解秩，$\lambda_r$为第$r$个成分的权重，$a_r$、$b_r$、$c_r$分别表示第$r$个成分在股票、因子和时间三个模态上的载荷向量。若将这些向量按列组合，可得因子矩阵$A=[a_1,\ldots,a_R]$、$B=[b_1,\ldots,b_R]$和$C=[c_1,\ldots,c_R]$。从经济含义上看，矩阵$A$描述股票在潜在空间中的位置，矩阵$B$反映原始因子对潜在成分的贡献强弱，矩阵$C$则揭示潜在结构随时间的变化轨迹。

CP 分解对应的目标可以写为
\begin{equation}
\min_{A,B,C,\lambda}\left\|\mathcal{X}-\sum_{r=1}^{R}\lambda_r a_r\circ b_r\circ c_r\right\|_F^2.
\end{equation}
这类模型具备明显优势，整体结构简单直观，参数数量也容易把控，并且在满足唯一性相关条件时，所得分解结果也具备良好的实际解读价值\textsuperscript{[4][5]}。在具体算法实现过程中，本文借助交替最小二乘的思路，依照不同模态依次完成因子矩阵的更新工作。以固定$B$和$C$再对$A$进行更新计算为例，其表达式可写为
\begin{equation}
A\leftarrow X_{(1)}(C\odot B)\left[(C^{\top}C)\ast(B^{\top}B)\right]^{-1},
\end{equation}
其中$\odot$表示 Khatri--Rao 积，$\ast$表示 Hadamard 积。由此可以看出，CP 分解实质上是通过模态展开与矩阵优化，逐步逼近原始张量的低秩结构。

\subsection{Tucker 分解}

与 CP 分解不同，Tucker 分解将原始张量表示为核张量与多个因子矩阵的$n$模乘积，即
\begin{equation}
\mathcal{X}\approx \mathcal{G}\times_1 A\times_2 B\times_3 C.
\end{equation}
其中，$\mathcal{G}\in\mathbb{R}^{R_1\times R_2\times R_3}$为核张量，$A\in\mathbb{R}^{N\times R_1}$、$B\in\mathbb{R}^{F\times R_2}$、$C\in\mathbb{R}^{T\times R_3}$分别对应三个模态上的低维表示矩阵。核张量$\mathcal{G}$不再要求对角结构，因此能够刻画不同潜在成分之间更复杂的高阶交互关系。

Tucker 分解对应的优化目标可写为
\begin{equation}
\min_{\mathcal{G},A,B,C}\left\|\mathcal{X}-\mathcal{G}\times_1 A\times_2 B\times_3 C\right\|_F^2.
\end{equation}
该模型最为突出的特点，就是能够给不同模态设定互不相同的秩$(R_1,R_2,R_3)$。比如可以在股票模态中保留较高维度，用来体现市场横截面存在的各类差异，同时在因子模态与时间模态中选用更低维度，以此提取出整体的核心主导结构。借助这类多线性秩约束方式，Tucker分解在处理具备非对称特征、关联耦合性较强的金融数据时，拥有更强的适配能力。在实际求解过程中，一般先通过HOSVD或者截断奇异值分解得到初始因子矩阵，之后再采用交替迭代的方式，不断调整更新核张量以及各个模态对应的因子矩阵。

\subsection{两类模型的比较与本文选择}

从结构层面来讲，CP分解可以看作是Tucker分解里核张量呈现超对角结构时形成的特殊情形，二者同属于低秩张量表达体系，但是各自侧重的建模能力有所区别。CP分解更加倾向于借助少量共有成分解释整体数据结构，适合用来挖掘市场中相对稳定统一的潜在规律；Tucker分解则能够依靠核张量留存各类成分之间的交互关联，更便于刻画股票、因子与时间三者之间非对称的复杂联动关系。

结合本文设定的研究目标来看，CP分解可以构建出结构简单清晰、方便开展对比分析的低秩基准模型，而Tucker分解更加适用于开展主体实证分析工作。这是由于股票因子数据当中，一般同时存在行业走势差异、投资风格轮换以及市场阶段转变等多种层次的变化特征，单纯依靠秩一张量进行叠加组合，很难完整刻画数据当中蕴藏的全部结构特征。基于这一原因，后文研究将会同步呈现CP分解与Tucker分解的实证结果，同时将PCA方法当作二维数据降维的对比参照，分别从数据重构效果、因子排序有效性以及结构经济解释能力三个层面，对比区分各类研究方法各自的适用范围与使用局限。

\section{金融数据的张量表示方法}

开展股票因子相关研究时，最为贴合实际研究场景的张量构建形式，就是将股票设定为第一模态，因子设定为第二模态，时间设定为第三模态。我们将其中每一个元素$x_{ijt}$定义为第$i$只股票在第$t$个交易日所对应的第$j$个因子具体数值，也就是
\begin{equation}
\mathcal{X}=(x_{ijt})\in\mathbb{R}^{N\times F\times T}.
\end{equation}
这种数据表达形式具备三方面实用优势。首先，该形式能够留存不同股票在横截面上的对应关联，方便直观查看行业类型、投资风格以及市值规模等维度下的结构相近程度。其次，它可以保留各个因子之间的联动关系，便于筛选出冗余因子，也能找出存在相互替代作用的因子组合。最后，该形式还能留存时间维度上的数据连贯性与阶段变化特征，用来分析因子整体结构在牛熊行情切换、流动性波动以及政策调整过程中发生的各类变化。

如果再结合正式实证研究里划定的样本选取范围，我们也能够在统一的数据架构当中，搭建出沪深300、上证50以及中证500这三类样本池对应的因子张量，只需要在筛选研究样本时，参照指数过往成分股名单来划定对应的股票范围即可。本文的研究目的并不是搭建面向全市场的实时交易模型，而是完成兼顾数据原有结构的降维处理以及挖掘内在运行规律，所以研究过程中会更加注重保证样本统计标准统一、选取的时间区间完整、搭建的特征指标稳定，不会将不同样本池的数据混杂在一起开展建模分析。

\section{模型评价指标}

为了合理评判数据降维的实际效果，本文选用重构误差、解释方差、排序相关、信息比率以及滚动稳定性作为评价指标。其中重构误差主要用来衡量分解后的张量和原始张量二者之间的偏差大小，具体可以写为
\begin{equation}
\mathrm{MSE}=\frac{1}{NFT}\|\mathcal{X}-\hat{\mathcal{X}}\|_F^2.
\end{equation}
与之对应，解释方差反映模型对原始张量波动信息的保留程度，数值越大通常表示降维后保留的有效结构越多。

除此之外，本文除了考量模型的数据重构能力之外，还十分看重低维数据表达具备的排序实用价值，因此研究还选用Rank IC和IR两类指标，用来分析降维后得到的潜在评分和个股下期收益之间的关联程度。其中Rank IC以等级相关作为评判依据，比较贴合金融研究里侧重标的排序、而非精准预测具体数值的研究场景。如果某个模型拥有较高的方差解释水平，但是Rank IC数值偏低，就说明该模型更擅长还原原始数据数值，却不一定能够很好地适用于股票筛选排序相关研究。

最后，本文还引入滚动稳定性指标，用来衡量模型在不同时间区间内提取潜在结构的一致程度。在因子降维的研究过程中，倘若模型在各个时间窗口内得出的潜在结构差异过大，即便单次拟合的效果表现较好，也很难总结出能够重复沿用的研究结论。所以重构误差、解释方差、Rank IC和滚动稳定性一同组成了本文的综合评价体系，其中前两项指标侧重评判数据数值拟合效果，后两项指标则偏向金融实际含义解读与整体研究结果的稳健性。

\chapter{数据预处理与张量表示}

\section{数据来源与样本说明}

本文研究所用到的数据，主要来源于自行搭建的A股市场研究数据库。研究所需原始数据大多依托Baostock平台获取，之后再将这些基础数据整理划分成日度K线行情、企业财务资料、上市公司业绩公告、指数历史成分股信息、股票除权复权相关数据以及衍生因子数据表等多种数据形式。本文选定的实证研究样本池分别为HS300、SZ50和ZZ500。

从整体研究规划来看，本文计划将整体研究时间范围设定在2015年1月至2026年4月，以此全面探究不同的市场运行阶段；而在本文正文当中用来开展对比分析的核心短期实验样本，选取的时间区间为2026年3月2日到2026年3月30日。本文依旧将这三类样本池作为主要研究对象，是因为它们有着界限明确的投资风格、充足的市场流动性，同时指数成分股的划定规则也相对稳定，能够更加清晰地探究因子内部结构和不同样本池投资风格之间存在的对应关系。

从数据覆盖角度看，当前研究数据已经完成前复权 K 线、财务表、报告表和指数成分历史的整理，但原始层数据的可用跨度与当前正文重点讨论的短窗口实验之间仍然存在差异。也就是说，数据基础具备继续向长期窗口扩展的条件，而本文主要展示的比较结果仍集中于 2026 年 3 月的短窗口样本。这也是本文在局限性分析中专门强调“实验窗口口径仍需进一步统一”的原因。

需要着重说明的是，完成原始数据采集工作，和整理出可以直接投入模型训练的标准数据，二者并不是同一个概念。现阶段搭建训练张量依旧以核心主线因子面板数据为主体，而本文的数据拓展方案，除了前期用到的一阶市场代理变量、财务时点指标、业绩快报以及业绩预告相关事件特征之外，还进一步纳入了外部市场利率、月度宏观经济指标、个股分红数据、上市公司重大事项公告以及公告标题文本信息。为了避免模型训练过程中出现数据披露滞后造成偏差，同时杜绝未来信息提前泄露的问题，本文单独梳理了各类拓展特征的实际披露时间与正式可用时间，并将这些特征划分成市场代理指标、外部宏观数据、财务时点特征、企业分红行为以及公告事件文本这五大类别。所有特征都会按照\texttt{pub\_date}与\texttt{available\_date}的对应匹配规则纳入训练样本。这样的处理方式，既能够说明当前模型输入数据范围的形成原因，也能为后续继续引入更多精细化原始数据表提供合理依据。

在样本数据整理划分方面，本文没有将全部股票数据整合至同一个张量中进行统一训练，而是分别针对HS300、SZ50以及ZZ500单独搭建对应的因子张量。采用这种处理方式的原因是，不同指数样本池在市值规模特点、行业组成结构以及投资风格取向方面都有着十分明显的区别。如果直接将各类样本数据混合在一起训练，模型得到的潜在特征会优先捕捉样本池之间的差别，难以挖掘出股票因子本身内在的结构联系。而按照不同样本池分开建模，更便于在特征相似度更高的横截面数据里，分析张量分解方法的实际运用效果。

\section{股票池与时间窗口设定}

各类指数对应的样本范围并不是固定不变的，因此本文依托指数成分股历史变动数据来划定股票研究范围，不使用单一时间节点的静态名单。对于任意一个交易日，仅将当日纳入目标指数成分行列的股票划入研究样本当中。这种处理方式存在两点优势，第一能够避免借助当前成分股名单回溯历史数据，从而杜绝未来信息提前泄露的问题；第二也可以更加真实地体现出长期研究区间内指数样本池内部成分发生的实际变动情况。

从整体研究规划层面来说，最为合适的完整研究区间应当设定在2015年1月至2026年4月，借此全面观测多个不同市场发展阶段的变化特征。不过本文正文里开展的核心基准对比分析，依旧选用2026年3月的短期时间区间所得实验结果作为主要依据。所以后文最先阐述的各类实证研究结论，都需要视作短期样本区间内得出的研究结果，不能直接当作长期完整时间区间下的通用结论来使用。

\section{数据预处理}

开展数据预处理工作，主要目的是在完好保留股票、因子、时间三维原始数据结构的基础上，把原始因子面板整理成能够直接接入张量分解模型的规范稳定数据。和只进行单次静态数据清洗的方式不同，本文严格遵循标准化处理流程，先统一确定样本选取范围，再依次完成缺失值补齐、异常值调整与数据标准化操作，最后按照时间先后顺序划分出训练样本、验证样本和测试样本。表\ref{tab:preprocess-panel-summary}展示了三类正式样本池在输入模型前的数据面板规模，从中能够看出，三个样本池均已处理完毕，形成了没有数据缺失的完整短期窗口因子面板。

\begin{table}[htbp]
\centering
\caption{正式样本池因子面板规模与完整性}
\label{tab:preprocess-panel-summary}
\small
\begin{tabular}{lccccc}
\toprule
样本池 & 交易日数 & 股票数 & 因子数 & 因子观测值总数 & 缺失率 \\
\midrule
HS300 & 21 & 72 & 7 & 10\,584 & 0.00\% \\
SZ50  & 21 & 50 & 7 & 7\,350  & 0.00\% \\
ZZ500 & 21 & 500 & 7 & 73\,500 & 0.00\% \\
\bottomrule
\end{tabular}
\end{table}

从数据处理的先后顺序来看，本文所使用的预处理流程可参考图\ref{fig:preprocess-flow}。整套流程首先统一样本统计标准并与交易日期完成对齐匹配，再依次完成缺失值填充、异常值修正以及数据标准化处理，最后生成规范的三阶张量数据。采用固定流程分步处理，能够避免各类预处理操作交叉混杂，进而引发数据统计口径混乱的问题。

\begin{figure}[htbp]
\centering
\begin{tikzpicture}[x=1cm,y=1cm]
  \tikzset{
    flowbox/.style={
      draw,
      rounded corners,
      align=center,
      text width=2.85cm,
      minimum height=1.15cm,
      inner xsep=0.08cm,
      inner ysep=0.12cm
    }
  }
  \node[flowbox] (a) at (0,0) {原始因子面板\\统一股票池\\与交易日};
  \node[flowbox] (b) at (3.75,0) {缺失值处理\\前向填充\\后向填充\\行业中位数};
  \node[flowbox] (c) at (7.5,0) {异常值收缩\\截尾处理\\MAD 稳健收缩};
  \node[flowbox] (d) at (1.9,-2.4) {训练集标准化\\验证集与测试集\\复用训练参数};
  \node[flowbox] (e) at (5.65,-2.4) {按时间顺序切分\\训练/验证/测试\\并构造三阶张量};
  \draw[->, thick, shorten >=0.12cm, shorten <=0.12cm] (a.east) -- (b.west);
  \draw[->, thick, shorten >=0.12cm, shorten <=0.12cm] (b.east) -- (c.west);
  \draw[->, thick, shorten >=0.12cm, shorten <=0.12cm] (c.south) |- (e.east);
  \draw[->, thick, shorten >=0.12cm, shorten <=0.12cm] (d.east) -- (e.west);
\end{tikzpicture}
\caption{正式实验中的因子预处理流程}
\label{fig:preprocess-flow}
\end{figure}

首先进行缺失值处理时，股票因子出现数据空缺，大多是个股停牌、上市时间较短、财务数据未及时披露以及数据源覆盖不全等原因导致。倘若直接剔除存在缺失值的样本，容易造成时间窗口断裂，还会让各个样本池的数据分布失衡，因此本文制定了分层填补规则。设定原始因子数值为$x_{ijt}$，经过填补处理后的数值$\tilde{x}_{ijt}$可写为
\begin{equation}
\tilde{x}_{ijt}=
\begin{cases}
x_{ijt}, & x_{ijt}\text{ 已观测};\\
x_{ij,t^-}, & x_{ijt}\text{ 缺失且存在最近一期历史值};\\
x_{ij,t^+}, & x_{ijt}\text{ 缺失且不存在历史值但存在后续值};\\
\operatorname{med}_{g(i),j,t}, & \text{以上条件均不满足},
\end{cases}
\end{equation}
其中$g(i)$代表第$i$只股票所属的行业类别，$\operatorname{med}_{g(i),j,t}$指代在$t$时刻下该行业内因子$j$的截面中位数。这套填补规则既兼顾了数据在时间维度上的连续特征，也保障了同行业个股之间的数据具备合理的对比参考性。

其次进入异常值处理环节，估值类、换手率类以及未来收益类因子，很容易受到市场极端行情或是低基数效应的干扰。为了防止张量分解的整体结果被少数极端数值所影响，本文以训练样本集中的统计指标作为参照标准，对各类因子开展稳健缩尾处理。设定$\operatorname{med}_j$和$\operatorname{MAD}_j$分别代表因子$j$在训练集内的中位数与中位数绝对偏差，经过收缩调整后的因子数值$\hat{x}_{ijt}$表示为
\begin{equation}
\hat{x}_{ijt}
=\min\left\{\max\left(\tilde{x}_{ijt},\operatorname{med}_j-3\operatorname{MAD}_j\right),\operatorname{med}_j+3\operatorname{MAD}_j\right\}.
\end{equation}
这一处理能够在保留排序信息的同时抑制异常跳点对 Frobenius 范数的放大效应。

再者进行标准化处理时，本文只利用训练集求解位置参数与尺度参数，再将这组统一参数统一应用到验证集和测试集当中，以此规避未来信息泄露的问题。设$\mu_j^{\mathrm{train}}$和$\sigma_j^{\mathrm{train}}$分别为因子$j$在训练集中的均值与标准差，标准化之后得到的结果$z_{ijt}$定义为
\begin{equation}
z_{ijt}=\frac{\hat{x}_{ijt}-\mu_j^{\mathrm{train}}}{\sigma_j^{\mathrm{train}}+\varepsilon},
\end{equation}
其中$\varepsilon$是为避免分母取值为零所设置的极小常数。经过这一步处理后，不同量纲大小的因子都会被转换至可以相互对比的统一尺度，以此防止部分数值偏大的因子在分解运算中占据过高权重。

最后开展样本划分与数据边界界定工作时，本文并未采取全量样本统一拟合的处理方式，而是清晰划分出训练区间、验证区间与测试区间。其中训练集用于确定数据预处理规则以及拟合模型相关参数，验证集用来筛选确定最优分解秩，测试集仅作为独立外集开展最终效果评估。对于各类拓展特征而言，只有已经明确理清披露时间与实际可用交易时间对应规则的数据，才能够纳入训练输入数据当中，其余时间界定尚未成熟的补充信息暂时不整合至核心张量内。借助以上一系列规范设定，本文最大程度保证了实验执行标准、数据选取范围与模型评价体系三者之间保持高度一致。

\section{张量表示}

在全部预处理工作完成之后，本文把各个样本池对应的数据集整理为三阶张量$\mathcal{X}\in\mathbb{R}^{N\times F\times T}$。其中$N$代表该样本池依据对应交易日有效成分股名单统计得出的个股数量，$F$为因子面板内包含的因子总个数，$T$则代表时间维度下的交易天数。由于指数成分股会随时间发生变动，不同时间段内纳入统计的股票名单存在差异，因此在实际程序搭建过程中，本文通过统一编码搭配缺失值补位的处理手段，确保张量各个维度的结构形态始终保持固定不变。

这种张量数据表达形式和传统矩阵形式最核心的区别在于，时间维度不再单纯归属于样本编号范畴，而是和股票维度、因子维度一同构成挖掘潜在数据结构的核心维度。依托该形式得到的模型运算结果，除了能够得到重构误差这类数值型指标之外，还可以输出三个维度各自对应的潜在特征矩阵，能够直接为后续开展个股相似性研判、因子重要度判定以及市场时间状态演变分析提供可靠的数据支撑。

\chapter{基于张量分解的股票因子降维与模式发现方法}

\section{问题定义}

本文的研究问题可进行形式化定义：已知股票—因子—时间构成的三阶张量$\mathcal{X}$，在完整保留其多维结构信息的基础上，构建对应的低维表达形式$\hat{\mathcal{X}}$，使其不仅能够较好地还原原始张量数据，还可以生成在股票、因子以及时间层面均具备实际解释价值的潜在特征。和普通仅实现数据降维的研究问题相比，本文更加注重同时完成数据降维与内在规律挖掘这两项研究目标。

具体来说，本文主要完成三项细分研究任务。第一，在因子维度实现高效的数据压缩，以此削减高维因子组合中存在的数据冗余。第二，在股票维度挖掘内在相似结构，以此判断不同个股在隐性风格暴露层面的相近程度。第三，在时间维度划分市场运行阶段，探究潜在数据结构在长期样本区间内的演变趋势与突变特征。围绕这三项研究目标，传统二维PCA方法虽然能够搭建数据重构的基准结果，却无法对多维度数据做出合理的业务解释，因此本文引入张量分解方法开展相关研究。

\section{张量分解模型构建}

本文同步搭建了CP、Tucker以及PCA三种分析模型。其中CP模型是结构最为简易的张量分解基准模型，主要用于探究在强低秩约束条件下，三阶张量是否能够依靠少量公共成分实现有效表征。Tucker模型作为本文的核心研究方法，可在不同数据维度分别设置对应秩值，以此捕捉维度之间更为复杂的交互关联结构。PCA模型没有基于张量结构开展建模，但其在传统降维研究中具备出色的实际效果，因此本文将其保留，作为二维数据压缩的对照基准模型。

对 CP 分解而言，模型秩$R$越大，理论上可表示的结构越丰富，但解释性和压缩率会下降；秩过小则可能无法表示足够的结构信息。对 Tucker 分解而言，需要同时设置股票模态秩、因子模态秩和时间模态秩。本文在实验中并不依赖人工经验一次性固定秩，而是先给出候选 rank 集合，再在验证集上比较不同 rank 对应的 held-out MSE，选出表现最优的参数组合。这种做法能在一定程度上避免“为了论文结论而人为挑选 rank”的主观性问题。

在因子筛选与模式解释上，本文主要依赖因子模态矩阵$B$及其载荷大小。若某个潜在成分在$B$中对质量因子、估值因子和动量因子分别给出不同权重，则说明不同金融含义的因子被组合进不同潜在方向中。进一步地，可定义因子$j$的综合重要性为
\begin{equation}
w_j=\frac{\|b_{j:}\|_2}{\sum_{k=1}^{F}\|b_{k:}\|_2},
\end{equation}
并按$w_j$排序识别核心因子。这一做法虽然不直接等价于经济学中的因子定价显著性检验，但对于“哪些因子在低维潜在结构中承担了主要信息功能”这一问题具有直观解释力。

\section{模型求解与实现}

本文整体实验流程可归纳为五个步骤。第一步，读取各样本池对应的因子面板数据与指数成分历史数据，搭建得到原始三阶张量数据。第二步，严格按照预先划定的训练集、验证集、测试集划分规则，完成数据预处理与样本分割工作。第三步，分别对各类模型在备选秩值范围内开展训练与验证操作，筛选确定最优模型参数。第四步，整合训练集与验证集数据重新拟合确定最终模型，仅采用测试集数据开展独立外样本效果评估。第五步，统一输出模型各项评价指标、因子重要度结果、时间维度市场状态、备选股票排序结果以及各类可视化图表成果。

在算法实现层面，CP模型与Tucker模型的运算过程均依托迭代更新的方式完成。为了降低局部最优问题对模型最终结果产生的干扰，本文采用固定随机种子、标准化预处理协议的处理方式，同时完整留存模型的秩选择过程与评价结果记录。PCA模型作为二维基准方法，采用与前两类模型完全一致的测试集标准开展评估工作，以此保证三种方法的对比结果具备公平性。需要特别说明的是，本文的研究核心聚焦于学术研究型实证分析框架，而非面向实际应用的高频交易系统，因此在方法实现过程中，更注重保障结果的可复现性、流程的可追溯性以及方法对比的一致性，而非追求算法的极限运行速度。

\chapter{实验结果与分析}

\section{实验设置}

本次正式实验所选用的样本池包含沪深300、深证50以及中证500。本文正文首先选取2026年3月2日至2026年3月30日这段时间区间内的实验数据展开实证结果对比。本次研究中所有样本池均依托规范的因子面板数据完成张量构建，同时在统一的数据预处理、样本划分以及结果评价标准下，对CP、Tucker和PCA三种方法进行横向对比。其中CP模型与Tucker模型主要用于完成张量多维结构的建模分析，PCA模型则作为二维数据压缩的基准模型，用来对比各类模型的数据拟合效果。

在投资组合效果评价的统一标准上，本文严格遵循A股市场T+1的交易规则：依托$t$日能够获取的因子数据生成个股排序交易信号，在$t+1$交易日完成仓位构建，同时设定五个交易日的持仓周期来测算远期收益水平；交易成本按照单边万分之一的标准，结合实际换手率进行对应扣除。该设定方式不仅能够规避当日生成交易信号并立即交易所带来的前视偏差问题，还能让组合层面的实证研究结果更加贴合A股市场真实的实际交易约束条件。

从输出结果看，实验不仅生成数值指标，还形成了解释方差对比图、Rank IC 对比图、模型指标总览图、时间状态变化图、因子重要性热力图，以及组合累计净值对比图和回撤曲线对比图。这样的图文并行呈现方式有助于把论文从“结果表格堆砌”转向“结构模式分析”，也是本文更强调的一点。

\section{分解质量分析}

首先从数值重构能力层面展开分析，三大样本池得出的实验结果呈现出较为一致的规律。在沪深300样本池中，CP、Tucker、PCA三种方法的均方误差依次为0.9140、0.3406和趋近于0，对应的解释方差分别为0.0017、0.6280与1.0000；在深证50样本池中，三类方法的均方误差分别为0.9898、0.3160和趋近于0，对应解释方差为-0.0527、0.6639、1.0000；在中证500样本池中，三者均方误差依次是1.0091、0.4832和趋近于0，对应解释方差为-0.0541、0.4953、1.0000。由此能够明显看出，在本次短期时间窗口的实验结果里，PCA模型在三组样本池内均展现出最优的数据重构效果，而在两种张量分解方法当中，Tucker模型的表现始终稳定优于CP模型。

从上述实验结果能够得出，倘若仅将数值重构效果作为评判标准，二维数据压缩方法依旧具备十分突出的优势。究其原因，PCA算法在完成矩阵展开处理之后，会直接沿着数据方差最大的方向构建模型，本身就更容易实现最优的数值拟合效果；而张量类方法在坚守多维度结构约束的基础上，还需要协调平衡不同维度之间的数据表征能力，这也就导致其单纯的数据重构性能无法占据上风。不过这一结论并不能说明张量研究方法不具备实际研究价值。与之相反，当研究重心偏向数据结构解读与内在规律挖掘时，适当舍弃一部分重构精度，以此换取逻辑更加明晰的多维度潜在特征表达，是一种十分合理的研究思路。

\begin{figure}[!t]
\centering
\begin{tikzpicture}
\begin{groupplot}[
  group style={group size=3 by 1, horizontal sep=1.3cm},
  width=0.28\textwidth,
  height=6.0cm,
  ybar,
  ymin=-0.1,
  ymax=1.08,
  symbolic x coords={CP,Tucker,PCA},
  xtick=data,
  enlarge y limits={upper=0.06,lower=0.04},
  extra y ticks={0},
  extra y tick labels={},
  extra y tick style={grid=major, major grid style={black!65, line width=0.35pt}},
  point meta=y,
  nodes near coords={\pgfmathprintnumber[fixed,precision=2]{\pgfplotspointmeta}},
  nodes near coords,
  nodes near coords align={vertical},
  every node near coord/.append style={font=\tiny},
  tick label style={font=\small},
  ylabel style={font=\small},
  title style={font=\small},
  legend style={font=\small, draw=none}
]
\nextgroupplot[title={HS300},ylabel={解释方差}]
\addplot[fill=blue!55] coordinates {(CP,0.0017) (Tucker,0.6280) (PCA,1.0000)};
\nextgroupplot[title={SZ50}]
\addplot[fill=orange!70] coordinates {(CP,-0.0527) (Tucker,0.6639) (PCA,1.0000)};
\nextgroupplot[title={ZZ500}]
\addplot[fill=teal!65] coordinates {(CP,-0.0541) (Tucker,0.4953) (PCA,1.0000)};
\end{groupplot}
\end{tikzpicture}
\caption{三个正式样本池上的模型解释方差对比图}
\label{fig:variance}
\end{figure}

图~\ref{fig:variance} 更直观地展示了解释方差差异。PCA在三个不同样本池中的表现均占据明显优势；Tucker模型始终能够维持较高的正向解释方差，而CP模型仅能在沪深300样本中将解释方差维持在趋近于零的水平，在深证50与中证500样本中其解释方差依旧为负数。这一现象充分说明，在本次短期研究窗口的统计标准下，采用低秩结构的CP模型无法稳定刻画股票因子的内在关联结构，而Tucker模型凭借更加灵活的多线性秩约束条件，能够更为平稳地提取出数据内部的有效结构信息。

\section{模式发现结果分析}

除去数据数值层面的重构效果之外，本文更加重视由潜在表征所挖掘出的内在数据规律。结合因子重要性的输出结果可以发现，价值因子、质量因子以及波动率因子，在多个样本池与不同模型的测算结果中均长期排在前列，这也证明上述几类因子是当前所用因子体系当中的核心信息主体。以沪深300样本池对应的Tucker模型实证结果为例，波动率因子、价值因子与质量因子的平均重要程度依次位居前列，而动量因子在本次短期研究区间内的重要性基本趋近于零，由此能够看出该时间段内的市场运行特征，更多依托估值、基本面质量与价格波动相关信息，并不偏向中长期趋势类行情信息。

从股票维度的研究角度来看，借助潜在特征表达能够构建出不同个股之间的相似关联关系。虽然本文并未将股票相似性矩阵设定为最终的实验评价指标，但是通过查阅实验运行记录以及输出的相似股票组合结果能够发现，属于同一行业或是投资风格相近的股票，往往会在潜在特征空间当中呈现出相似度较高的载荷分布结构。这一研究结论，也和量化投资领域内“因子暴露情况决定股票风格归类”的基本研究思路保持一致。与传统简单相关系数的衡量方式相比，依托张量潜在表征所构建的股票相似关系，能够同时兼顾时间维度与因子维度两大层面的信息，因此可以更加精准地体现出股票之间深层次的结构相近程度，而不只是单纯反映短期内股票价格走势的同步变化情况。

依托全A股样本开展Tucker模型实证分析后，能够绘制出股票潜在结构图以及聚类与行业交叉分布图。其中前者按照股票的潜在聚类标签进行排布，同时借助不同行业对应的色彩，直观展现各类聚类内部的行业组成情况；后者则针对聚类标签与行业分布情况开展交叉统计分析。这两类可视化结果共同证实，股票在潜在特征空间中的划分结果，并不会单纯复刻传统行业分类标准，而是在行业信息的基础上，进一步保留了由因子暴露程度、市场时间运行状态以及个股排序得分共同形成的内在结构相似特征。

\begin{figure}[!t]
\centering
\begin{tikzpicture}
\begin{groupplot}[
  group style={group size=2 by 1, horizontal sep=2.2cm},
  width=0.36\textwidth,
  height=7.0cm,
  xbar,
  enlarge y limits=0.18,
  yticklabel style={font=\scriptsize, text width=2.0cm, align=right},
  tick label style={font=\small},
  label style={font=\small},
  title style={font=\small},
  point meta=x,
  nodes near coords={\pgfmathprintnumber[fixed,precision=2]{\pgfplotspointmeta}},
  nodes near coords,
  every node near coord/.append style={font=\tiny},
]
\nextgroupplot[
  title={Tucker跨边界Rank IC},
  xmin=-0.42,
  xmax=0.28,
  symbolic y coords={大市值,小市值,电子C39,专用设备C35,中市值,中证500,全市场,沪深300,上证50,医药C27},
  ytick=data,
  xlabel={Rank IC}
]
\addplot[fill=blue!55] coordinates {
(-0.3780,大市值)
(-0.0997,小市值)
(-0.1196,电子C39)
(-0.0496,专用设备C35)
(0.1031,中市值)
(0.0209,中证500)
(0.0279,全市场)
(-0.0468,沪深300)
(0.0468,上证50)
(0.1476,医药C27)
};
\nextgroupplot[
  title={Tucker跨边界累计收益},
  xmin=-12,
  xmax=32,
  symbolic y coords={大市值,小市值,电子C39,中市值,全市场,中证500,上证50,专用设备C35,沪深300,医药C27},
  ytick=data,
  xlabel={累计收益(\%)}
]
\addplot[fill=orange!70] coordinates {
(-7.26,小市值)
(-3.65,电子C39)
(-0.28,大市值)
(-2.11,全市场)
(-1.63,中证500)
(-0.43,中市值)
(0.04,上证50)
(0.17,专用设备C35)
(1.29,沪深300)
(16.46,医药C27)
};
\end{groupplot}
\end{tikzpicture}
\caption{短窗口样本边界扩展下 Tucker 的跨边界排序与收益分化}
\label{fig:tucker_boundary_short}
\end{figure}

时间模态则提供了对市场状态切换的观察窗口。HS300 的 Tucker 结果显示，最明显的相邻日期结构跳变集中在 2026 年 3 月 25 日至 2026 年 3 月 30 日这一短区间，说明即使在较短的正式输出窗口内，潜在结构也并非静态稳定，而是在局部阶段出现明显重排。换言之，张量分解不仅用于“压缩因子”，还能够识别“市场何时发生了潜在结构变化”。

\begin{figure}[!t]
\centering
\begin{tikzpicture}
\begin{groupplot}[
  group style={group size=3 by 1, horizontal sep=1.3cm},
  width=0.28\textwidth,
  height=6.0cm,
  ybar,
  ymin=-0.10,
  ymax=0.18,
  symbolic x coords={CP,Tucker,PCA},
  xtick=data,
  enlarge y limits={upper=0.10,lower=0.08},
  enlarge x limits=0.18,
  point meta=y,
  nodes near coords={\pgfmathprintnumber[fixed,precision=2]{\pgfplotspointmeta}},
  nodes near coords,
  nodes near coords align={vertical},
  every node near coord/.append style={font=\tiny},
  tick label style={font=\small},
  ylabel style={font=\small},
  title style={font=\small}
]
\nextgroupplot[title={HS300},ylabel={Rank IC}]
\addplot[fill=blue!55] coordinates {(CP,-0.0278) (Tucker,-0.0468) (PCA,-0.0530)};
\nextgroupplot[title={SZ50}]
\addplot[fill=orange!70] coordinates {(CP,0.0198) (Tucker,0.0468) (PCA,-0.0224)};
\nextgroupplot[title={ZZ500}]
\addplot[fill=teal!65] coordinates {(CP,-0.1082) (Tucker,0.0209) (PCA,-0.0442)};
\end{groupplot}
\end{tikzpicture}
\caption{三个正式样本池上的模型 Rank IC 对比图}
\label{fig:rankic}
\end{figure}

图~\ref{fig:rankic} 展示了三个正式样本池在 Rank IC 指标上的差异结构。可以看到，在严格 T+1 口径下，HS300 的三类模型 Rank IC 均转为负值，但幅度仍集中在较低区间；SZ50 与 ZZ500 中则分别只有 Tucker 保持了 0.0468 和 0.0209 的正向排序相关。这意味着仅有较强重构能力并不足以保证排序有效性，保持三维结构的潜在表示在部分样本池中仍具有独立价值。

\section{选股有效性分析}

为从金融实际应用层面评判降维效果，本文进一步选用Rank IC指标开展量化检验。图~\ref{fig:rankic}结果表明，在沪深300样本内，CP、Tucker、PCA三种方法的Rank IC均值依次为$-0.0278$、$-0.0468$与$-0.0530$，这也体现出大盘蓝筹类股票在严格T+1交易规则下，短期区间内的收益排序难度相对更大；在深证50样本中，Tucker模型Rank IC数值为0.0468，CP模型为0.0198，而PCA模型则降至$-0.0224$；在中证500样本当中，也仅有Tucker模型能够维持0.0209的正向排序相关水平。由此能够看出，以跨不同样本池稳定保持正向排序相关性作为评判标准时，Tucker模型在本次实证区间内的综合表现最为均衡稳定。

该实验结果表明，数值拟合能力最优的方法，并不等同于股票收益排序效果最优的方法。PCA模型在解释方差层面具备明显优势，但其在中证500样本中的Rank IC依旧为负值，这说明该方法的数据压缩思路更偏向于对静态数据方差进行解释，难以直接应用于股票横截面收益排序研究。与之相对，Tucker模型依托完整的三维结构开展建模，在实现因子维度压缩的同时，保留了更多贴合时间演变规律与股票截面特征的内在规律，因此该模型在三类样本池内均未出现Rank IC由正向转为负向的现象。

站在量化选股的研究视角来看，该研究结论具备较强的现实实践意义。在实际的量化实证研究当中，我们并非要求数据降维得到的低维表征能够精准还原每一组原始因子数值，而是更加看重压缩后得到的潜在特征，能否支撑起稳定性更强的股票收益排序以及市场投资风格的合理解释。倘若某一模型仅能做到重构误差最小，选股排序效果却波动较大，那么该模型仅能算作单纯的数值压缩工具；而如果模型在选股排序指标与滚动实证稳定性方面都表现得更为均衡，就代表其提取出的潜在特征更适用于后续金融逻辑解读与股票截面投资决策。本次实证结果能够证实，在短期时间窗口的横向对比实验里，Tucker模型在这一应用层面具备更为突出的优势。

进一步把排序结果转化为组合层指标后，可以观察到“排序相关性”和“组合表现”并不完全同步。在基准方案结果中，HS300 上三类模型组合都取得了正累计收益，其中 Tucker 达到 1.29\%，高于 PCA 的 0.92\% 和 CP 的 0.53\%；SZ50 上只有 Tucker 保持微弱正收益 0.04\%，而 CP 与 PCA 分别回落到 $-2.27\%$ 和 $-0.56\%$；ZZ500 上 Tucker 与 PCA 的累计收益分别为 $-1.63\%$ 和 $-1.82\%$，虽然都未转正，但仍显著好于 CP 的 $-7.76\%$。换言之，单纯较高的 Rank IC 并不自动等价于更优的组合净值，这也是将组合结果与 Rank IC 联动分析的必要性所在。

引入拓展特征后的实验结果，进一步凸显了不同模型间的效果差异。在沪深300样本中，加入拓展特征后三类模型的组合收益表现均出现下滑，CP、Tucker、PCA的累计收益依次降至$-1.88\%$、$-3.12\%$和$-2.56\%$，对应的Rank IC指标也全部转为负值；在深证50样本中，新增输入特征同样未能优化选股排序效果，三类模型累计收益分别为$-3.07\%$、$-4.12\%$与$-3.49\%$；中证500样本的实验结果则出现明显分化，CP模型累计收益回升至$1.18\%$，Rank IC同步上升至0.0342，而Tucker与PCA模型收益依旧回落至$-3.39\%$和$-3.47\%$。这一结果证明，拓展特征不仅会改变模型的数据重构能力与选股排序效果，还会调整投资组合层面的收益与风险配比关系。故而在评判新增输入特征的实际作用时，需要综合考量净值走势、最大回撤、换手率以及行业、风格敞口等多项维度，不可只凭借单一排序指标完成效果判定。

在进一步补充分位数组、多空组合、成本后净值和基准超额收益之后，可以更细致地判断“排序有效性是否真正能够转化为可配置收益”。例如，基准方案下的 HS300 样本中，三类模型的 \texttt{top-bottom} 多空净值全部为正，其中 PCA、Tucker 与 CP 的多空累计净值分别达到 1.85\%、1.71\% 和 1.68\%；但一旦计入交易成本，三类模型的成本后累计收益分别回落到 0.91\%、1.28\% 和 0.51\%，说明毛收益与换手约束并不总是同向变化。又如基准方案下的 SZ50 样本中，Tucker 的多空组合累计净值达到 7.50\%，成本后净值仍基本持平，并保留了 0.04\% 的微弱基准超额，说明其分位排序结构虽弱但并未被基准完全吞没。扩展方案下的 ZZ500 样本中，Tucker 与 PCA 的成本后净值分别回落到 $-3.40\%$ 和 $-3.48\%$，相对基准的超额净值也分别为 $-5.40\%$ 和 $-5.48\%$，说明更宽输入边界虽然改变了结构解释，但在当前短窗口中并未把中盘样本稳定转化为可验证的正向配置价值。可见，加入分位数组、多空、成本和超额收益之后，组合层证据已经比单纯的前N净值曲线更接近真实投资决策场景。

为了把组合层结果与样本边界、行业暴露和风格暴露放到同一视角下比较，本文进一步构造了统一的跨边界汇总表。该表同时列出各样本池的 Rank IC 最优模型、累计收益最优模型、Sharpe 最优模型，以及每个模型的前N组合收益、分组价差、多空净值、成本后净值、平均交易成本、超额净值、年化波动、平均换手、最大回撤和主要行业/风格暴露。汇总结果显示，Tucker 仍在 SZ50、ZZ500 和医药行业样本中兼具较好的排序和收益表现，但在 HS300 上最优累计收益与 Sharpe 来自 Tucker、最优 Rank IC 则来自 CP；在全 A 样本中收益与 Sharpe 最优为 PCA；在大市值与小市值样本中，Rank IC 和收益最优也再次发生分离。这说明模型优劣并不是单一指标可以决定的，必须同时观察排序相关性、组合收益、换手成本、超额表现和暴露集中度。

将本文研究成果与已有相关文献进行对照分析，能够更加清晰地凸显本研究的补充价值。黄宏运与刘宇轩的相关研究证实，即便在传统回归分析框架下搭建多因子模型，依靠因子筛选搭配宏观周期修正的手段，依旧可以取得不错的选股实证效果\textsuperscript{[13][14]}。本文研究结论并不否认这一观点，而是进一步提出，当研究对象从经过筛选的少量优质因子拓展至高维因子集合后，研究核心便从验证因子有效性，转变为探究如何对高维信息完成稳健压缩，并将其应用于股票收益排序研究\textsuperscript{[15][16]}。与此同时，岳欣与曾亚丽围绕金融数据张量化展开的研究，着重强调了时间结构与高阶表征的重要研究价值\textsuperscript{[8][10]}。而本文在拓展研究样本范围、开展长窗口年度复核以及分析时间状态切换规律中得到的结论进一步说明，高阶结构带来的研究优势，不仅能够体现在市场波动分析、时序预测领域，同样也可有效应用于股票因子排序与市场模式识别相关研究当中\textsuperscript{[9][11][12][17]}。简言之，本文致力于在多因子量化研究与金融数据张量表征研究之间，建立更为直观紧密的实证关联。

\section{扩展特征对照分析}

为探究拓宽信息覆盖范围是否会对模型运行效果产生影响，本文进一步构建拓展因子面板，分别在沪深300、深证500、中证500三个样本池内开展基准实验与拓展对照实验。拓展实验在原有四类核心因子的基础上，新增市场状态代理变量、外部宏观经济指标、时点财务指标、分红相关指标，同时纳入业绩快报、业绩预告以及公告事件文本特征等各类信息。从具体数据维度来看，本次新增输入指标涵盖短中期收益动量、股价波动率、资产回撤水平、成交额变动、政策利率、物价与货币供给相关指标，还包含企业盈利水平、上市公司分红特征以及公告事件活跃度等多个研究层面。两组实验的对比结果可以发现，引入拓展信息无法单方面优化模型效果，其产生的作用在不同样本池当中呈现出十分显著的差异化特征。

从已生成的对照结果看，HS300 上新增市场代理变量、外部宏观、分红和公告事件字典后，三类模型的 Rank IC 全部转负，CP、Tucker 与 PCA 分别回落到 $-0.2851$、$-0.1417$ 和 $-0.1073$，组合累计收益也同步转弱，其中 CP、Tucker 与 PCA 分别为 $-1.88\%$、$-3.12\%$ 和 $-2.56\%$，说明扩展输入与大盘蓝筹样本之间仍存在显著相容性差异。SZ50 上，扩展输入并未改善排序表现，CP、Tucker 与 PCA 的 Rank IC 分别为 $-0.0345$、$-0.0509$ 与 $-0.1223$，累计收益也分别降至 $-3.07\%$、$-4.12\%$ 和 $-3.49\%$，表明“输入扩展”并没有对所有模型产生一致收益。ZZ500 上，扩展输入的分化更明显：CP 的 Rank IC 升至 0.0342，累计收益转为 1.18\%，而 PCA 与 Tucker 的 Rank IC 仍为负，累计收益分别为 $-3.47\%$ 和 $-3.39\%$。因此，扩展输入的作用更像是重新塑造不同模型在不同样本池中的信息利用方式，而不是简单地把所有结果整体抬升。

\begin{figure}[!t]
\centering
\begin{tikzpicture}
\begin{groupplot}[
  group style={group size=3 by 1, horizontal sep=1.0cm},
  width=0.29\textwidth,
  height=5.8cm,
  ybar,
  ymin=-0.20,
  ymax=0.12,
  symbolic x coords={CP,Tucker,PCA},
  xtick=data,
  enlarge y limits={upper=0.10,lower=0.08},
  enlarge x limits=0.18,
  point meta=y,
  nodes near coords={\pgfmathprintnumber[fixed,precision=2]{\pgfplotspointmeta}},
  xticklabel style={font=\tiny},
  ylabel style={font=\small},
  title style={font=\small},
  legend style={font=\tiny, draw=none, at={(0.5,1.08)}, anchor=south},
  nodes near coords,
  every node near coord/.append style={font=\tiny},
]
\nextgroupplot[title={HS300},ylabel={Rank IC}]
\addplot[fill=blue!55] coordinates {(CP,-0.0278) (Tucker,-0.0468) (PCA,-0.0530)};
\addplot[fill=orange!70] coordinates {(CP,-0.2851) (Tucker,-0.1417) (PCA,-0.1073)};
\legend{基准方案,扩展方案}
\nextgroupplot[title={SZ50}]
\addplot[fill=blue!55] coordinates {(CP,0.0198) (Tucker,0.0468) (PCA,-0.0224)};
\addplot[fill=orange!70] coordinates {(CP,-0.0345) (Tucker,-0.0509) (PCA,-0.1223)};
\nextgroupplot[title={ZZ500}]
\addplot[fill=blue!55] coordinates {(CP,-0.1082) (Tucker,0.0209) (PCA,-0.0442)};
\addplot[fill=orange!70] coordinates {(CP,0.0342) (Tucker,-0.0138) (PCA,-0.0189)};
\end{groupplot}
\end{tikzpicture}
\caption{三组正式样本池在基准方案与扩展方案下的 Rank IC 对比}
\label{fig:extended_rankic}
\end{figure}

\begin{figure}[!t]
\centering
\begin{tikzpicture}
\begin{groupplot}[
  group style={group size=3 by 1, horizontal sep=1.0cm},
  width=0.29\textwidth,
  height=5.8cm,
  ybar,
  ymin=-10.5,
  ymax=5.0,
  symbolic x coords={CP,Tucker,PCA},
  xtick=data,
  enlarge y limits={upper=0.10,lower=0.06},
  enlarge x limits=0.18,
  point meta=y,
  nodes near coords={\pgfmathprintnumber[fixed,precision=2]{\pgfplotspointmeta}},
  xticklabel style={font=\tiny},
  ylabel style={font=\small},
  title style={font=\small},
  legend style={font=\tiny, draw=none, at={(0.5,1.08)}, anchor=south},
  nodes near coords,
  every node near coord/.append style={font=\tiny},
]
\nextgroupplot[title={HS300},ylabel={累计收益(\%)}]
\addplot[fill=blue!55] coordinates {(CP,0.53) (Tucker,1.29) (PCA,0.92)};
\addplot[fill=orange!70] coordinates {(CP,-1.88) (Tucker,-3.12) (PCA,-2.56)};
\legend{基准方案,扩展方案}
\nextgroupplot[title={SZ50}]
\addplot[fill=blue!55] coordinates {(CP,-2.27) (Tucker,0.04) (PCA,-0.56)};
\addplot[fill=orange!70] coordinates {(CP,-3.07) (Tucker,-4.12) (PCA,-3.49)};
\nextgroupplot[title={ZZ500}]
\addplot[fill=blue!55] coordinates {(CP,-7.76) (Tucker,-1.63) (PCA,-1.82)};
\addplot[fill=orange!70] coordinates {(CP,1.18) (Tucker,-3.39) (PCA,-3.47)};
\end{groupplot}
\end{tikzpicture}
\caption{三组正式样本池在基准方案与扩展方案下的组合累计收益对比}
\label{fig:extended_return}
\end{figure}

图~\ref{fig:extended_rankic}与图~\ref{fig:extended_return}能够共同证实，引入拓展特征数据并未出现整体提升或是整体下滑的单一变化趋势，还进一步加大了不同模型与各类样本池之间的效果差异。这一现象说明，扩充数据信息维度确实能够改变数据内部的潜在结构，但这类结构上的调整能否切实提升模型整体性能，还会受到样本池市场风格、模型自身架构以及特征筛选方法等多重因素的制约。基于上述实验结论，本文将该组实证结果认定为拓宽了数据输入范围，但仍需进一步筛选有效特征指标的阶段性研究依据。

\section{样本边界扩展探索}

为进一步验证本文所得研究结论是否仅适用于单一指数样本范围，本文增设了样本边界拓展实验，实验样本涵盖全市场A股活跃个股样本、三类行业分层样本以及大中小市值分层样本。相较于沪深300、深证500、中证500所构成的基准研究样本，该组拓展实验着重探究模型在不同样本划分标准下运行效果的稳定性。研究先在短期时间统计口径下完成各组实验结果对比，再依托长期因子面板数据，对2015至2026年所有完整年度数据开展逐年复核分析，重点观测Rank IC数值、策略滚动运行稳定性以及投资组合累计收益的波动与相对变动情况。

从全A股样本的实证结果可以看出，Tucker模型依旧是三种方法里选股排序表现最为稳健的一种，其Rank IC数值达到0.0279，明显优于CP模型的$-0.0978$与PCA模型的$-0.0490$；在累计收益方面，该模型收益为$-2.11\%$，尽管略低于PCA模型的$-1.82\%$，但整体表现显著优于CP模型的$-5.52\%$。这一结果能够说明，当研究样本范围从指数成分股拓展至覆盖面更广的全市场A股活跃个股时，依托三维结构建模所具备的优势并未消失，这类优势更多体现在股票排序的稳定程度上，而非单纯实现组合收益的数值最优。

行业分层样本的实证结果呈现出更为明显的结构性差异。在电子行业样本（C39）中，三类模型的Rank IC指标均为负值，其中CP模型数值为$-0.0164$，表现明显优于Tucker模型的$-0.1196$与PCA模型的$-0.1292$，这也说明在高波动成长类行业中，个股截面选股信号容易受到特定模型结构的干扰偏离。在医药行业样本（C27）内，Tucker模型的Rank IC达到0.1476，累计收益高达$16.46\%$，整体表现远超CP模型与PCA模型，能够看出在盈利质量优势显著、事件驱动特征鲜明的行业当中，三维结构压缩与股票排序之间可以形成更强的协同作用。在专用设备行业样本（C35）中，CP模型取得最高Rank IC数值0.0520，而Tucker模型尽管Rank IC仅为$-0.0496$，依旧实现了$0.17\%$的小幅正向收益。综上可知，不同行业研究场景下不存在固定不变的最优模型，样本本身的风格特征会显著影响张量分解方法与二维压缩方法的相对优劣水平。

市值分层实验结果进一步证实，样本选取范围的差异，会改变选股排序效果与组合实际收益表现之间的对应关系。在小市值股票样本中，CP模型的Rank IC达到0.1399，但其组合累计收益仅为$-1.06\%$，这说明在市场噪声偏高的环境下，理想的排序相关度并不一定能够转化为可实际落地获取的投资收益。在中市值股票样本里，PCA模型的Rank IC为0.1224，而组合累计收益表现最佳的却是CP模型，收益数值为$1.46\%$，Tucker模型则只取得$-0.43\%$的小幅负收益，这意味着中盘股票样本更容易出现排序效果最优与实际收益最优相互背离的现象。在大市值股票样本当中，CP模型拥有最高的Rank IC，数值为0.1989，但三类模型搭建的投资组合收益均未能明显实现正向盈利。这一结果再次说明，单一排序指标表现突出，和投资组合最终净值实现正向收益之间，依旧存在十分明显的差距。

在全年度长期数据复核过程中，模型最优表现随样本范围与研究年份发生切换的现象依旧普遍存在。全A股样本在2015年区间内，Rank IC、组合收益与夏普比率三项指标的最优结果均由CP模型取得；2019年三类模型的累计收益全部转为正向收益，此时表现最优的模型变为PCA；2021年则迎来Tucker模型同时在Rank IC、收益以及夏普比率三项维度占据优势的阶段；而2026年短期年度区间内又形成指标分化格局，PCA在Rank IC指标上领先，Tucker在组合收益层面更具优势，CP则在夏普比率指标中表现更佳。不论是行业分层样本还是市值分层样本，都不存在能够在所有年份里始终保持最优表现的固定模型。这足以说明，拓展研究样本边界不仅会改变各类模型的选股排序效果，还会使得不同市场行情阶段下具备竞争优势的模型分布格局发生相应变化。

\begin{figure}[!t]
\centering
\begin{tikzpicture}
\begin{axis}[
  width=0.82\textwidth,
  height=7.0cm,
  ymin=0.7,
  ymax=2.4,
  symbolic x coords={2015,2019,2021,2026},
  xtick=data,
  xlabel={代表年份},
  ylabel={成本后累计净值},
  legend style={font=\small, draw=none, at={(0.5,1.05)}, anchor=south, legend columns=3},
  tick label style={font=\small},
  label style={font=\small},
  title style={font=\small},
  mark size=2.2pt,
]
\addplot+[mark=square*, blue, thick] coordinates {(2015,0.9871) (2019,1.4637) (2021,1.7381) (2026,0.9080)};
\addplot+[mark=triangle*, orange, thick] coordinates {(2015,0.8530) (2019,1.7245) (2021,1.9464) (2026,0.8600)};
\addplot+[mark=*, teal!70!black, thick] coordinates {(2015,0.6749) (2019,1.4228) (2021,2.1971) (2026,0.9194)};
\legend{CP,PCA,Tucker}
\end{axis}
\end{tikzpicture}
\caption{全 A 长窗口代表年份的成本后累计净值对比}
\label{fig:all_a_long_window_cost}
\end{figure}

图~\ref{fig:all_a_long_window_cost} 进一步给出了 2015、2019、2021 与 2026 四个代表年份下全 A 长窗口运行的成本后累计净值。该图最直接地展示了一个事实：模型优劣并不会在所有市场阶段保持不变。2015 年只有 CP 尚能把成本后净值维持在 1 附近，而 Tucker 明显落后；2019 年三类模型全部转正，但领先者改为 PCA；2021 年则出现 Tucker 明显跑赢 CP 与 PCA 的结构；到 2026 年短窗口时，三类模型又重新回落到 1 以下。换言之，长窗口复核真正补充的不是“某一模型永远最好”，而是“不同阶段的优势模型会发生切换”，这恰恰是样本边界扩展实验最重要的稳健性信息。

从因子暴露层面分析可知，行业分层样本的行业持仓敞口会自然集中在对应单一行业之中，而全A股样本与指数成分股样本，则更能清晰体现出不同模型的选股偏好与风格倾向。例如依托全A股数据构建的Tucker模型组合，行业持仓主要集中在医药制造、黑色金属、影视传媒等领域，风格暴露也以质量因子为主；沪深300样本下的Tucker组合则明显偏向货币金融服务行业，同时融合了质量、价值与波动率三类因子风格。

图~\ref{fig:tucker_boundary_short}进一步呈现了短窗口内拓展样本边界后，Tucker模型在不同场景下排序能力与组合收益的分化情况：医药行业样本能够同时在Rank IC和累计收益两项指标上占据优势，而电子行业、小市值以及大市值样本，即便模型依旧可以捕捉到部分市场结构信号，也难以将这类信号稳定转化为正向的组合净值收益。这一结果也进一步说明，开展市场模式挖掘研究时，不能只判定潜在聚类结构是否稳定，还需要核查投资组合是否会因样本边界变动，出现持仓扎堆少数行业或是局限于单一投资风格的问题。

\begin{figure}[!t]
\centering
\begin{tikzpicture}
\begin{groupplot}[
  group style={group size=2 by 1, horizontal sep=2.2cm},
  width=0.36\textwidth,
  height=6.6cm,
  xbar,
  enlarge y limits=0.18,
  yticklabel style={font=\scriptsize, text width=2.0cm, align=right},
  tick label style={font=\small},
  label style={font=\small},
  title style={font=\small},
  point meta=x,
  nodes near coords={\pgfmathprintnumber[fixed,precision=2]{\pgfplotspointmeta}},
  nodes near coords,
  every node near coord/.append style={font=\tiny},
]
\nextgroupplot[
  title={2026 长窗口 Tucker Rank IC},
  xmin=-0.16,
  xmax=0.14,
  symbolic y coords={电子C39,小市值,全市场,专用设备C35,大市值,医药C27,中市值},
  ytick=data,
  xlabel={Rank IC}
]
\addplot[fill=teal!65] coordinates {
(-0.1215,电子C39)
(-0.0877,小市值)
(-0.0708,全市场)
(-0.0462,专用设备C35)
(0.0102,大市值)
(0.0161,医药C27)
(0.0814,中市值)
};
\nextgroupplot[
  title={2026 长窗口 Tucker 期末净值},
  xmin=0.68,
  xmax=1.30,
  symbolic y coords={小市值,电子C39,中市值,专用设备C35,全市场,大市值,医药C27},
  ytick=data,
  xlabel={期末净值}
]
\addplot[fill=purple!45] coordinates {
(0.7940,小市值)
(0.8156,电子C39)
(0.8745,中市值)
(0.8807,专用设备C35)
(0.9196,全市场)
(0.9260,大市值)
(1.2818,医药C27)
};
\end{groupplot}
\end{tikzpicture}
\caption{2026 年长窗口结果中 Tucker 的跨边界表现}
\label{fig:tucker_boundary_long2026}
\end{figure}

\section{稳健性与局限性分析}

在滚动稳定性层面，沪深300、深证500与中证500样本中，Tucker模型稳定性数值依次为0.9898、0.9923和0.9683，整体水平显著优于CP模型，同时与PCA模型基本持平。模型具备良好稳定性，意味着其在不同滚动区间内挖掘得到的潜在数据结构具备高度一致性，能够让研究人员更为合理地解读因子重要程度与时序市场状态的演变规律，避免因模型在不同窗口内出现明显偏移，致使实证结果失去合理的解释依据。对于学术实证研究而言，模型运行的稳定特质远比单一区间的优异表现更具价值，也是保障相关研究结论具备可重复性的关键条件。

\begin{figure}[htbp]
\centering
\begin{tikzpicture}
\begin{groupplot}[
  group style={group size=3 by 1, horizontal sep=1.3cm},
  width=0.28\textwidth,
  height=6.0cm,
  ybar,
  ymin=0.50,
  ymax=1.02,
  symbolic x coords={CP,Tucker,PCA},
  xtick=data,
  nodes near coords,
  nodes near coords align={vertical},
  every node near coord/.append style={font=\tiny},
  tick label style={font=\small},
  ylabel style={font=\small},
  title style={font=\small}
]
\nextgroupplot[title={HS300},ylabel={滚动稳定性}]
\addplot[fill=blue!55] coordinates {(CP,0.6974) (Tucker,0.9898) (PCA,0.9946)};
\nextgroupplot[title={SZ50}]
\addplot[fill=orange!70] coordinates {(CP,0.5320) (Tucker,0.9923) (PCA,0.9956)};
\nextgroupplot[title={ZZ500}]
\addplot[fill=teal!65] coordinates {(CP,0.8384) (Tucker,0.9683) (PCA,0.9575)};
\end{groupplot}
\end{tikzpicture}
\caption{三个正式样本池上的模型滚动稳定性对比图}
\label{fig:stability}
\end{figure}

当然，本篇论文还有不少内容能够进行后续的优化调整，首先从数据录入的相关角度来说，研究人员在拓展实验的过程当中，已经将市场参考指标、外界利率数据、月度宏观相关数据、财务PIT数据、分红相关信息以及公告文本整理出的各类特征，全都加入到训练使用的数据当中，还制定好了清晰的信息公开时间与实际使用时间的相关限制条件，不过外界宏观相关资料和各类公告文字内容，大多只收集了本次实验时间段周边的相关内容，并没有整理收录更早阶段的历史数据，也没有集齐全部完整的相关文本资料，这样的情况也能看出，目前使用的研究模型，已经证实范围更广的录入数据，能够在时间节点相互契合的基础下，纳入到训练样本当中，而想要确定这套模型能否长期适配范围更广、来源各不相同的各类指标体系，还需要不断补充完善基础原始数据表，之后再开展对应的验证测试工作。

第二，在效果验证层面，本文已经补齐了前N组合的累计净值、回撤、换手率和行业/风格暴露等基础结果，并进一步加入了分位数组、多空组合、成本后净值、交易成本和相对样本池指数的超额收益，且已通过跨样本边界汇总表与 \texttt{2015-2026} 全年长窗口组合汇总统一比较。该证据链已经能够从短窗口与全年窗口两个层面回答“若将排序结果转化为真实组合规则，会得到怎样的收益—风险画像”；当前仍需继续深化的是更细粒度的交易冲击刻画和风险归因模型，而不是再停留在“是否具备长窗口证据”这一层面。

第三，在样本边界层面，本文已经从 HS300、SZ50 与 ZZ500 三个指数样本池扩展到全 A 活跃股票、行业分层和市值分层实验，并为 2015--2026 年区间完成了按年度拆分的长窗口复核。新增结果说明，指数样本上的结论不能被机械外推：Tucker 在部分窗口和部分行业样本中表现较稳健，但电子、专用设备和不同市值分层样本的最优模型并不固定。虽然全年长窗口结果已经补齐，但当前训练输入仍属于既定研究边界下的正式比较，因此这批扩展实验更适合被理解为“现有输入范围在不同边界与市场阶段下的稳健性检验”，而不是所有候选输入都已完备后的最终全市场结论。

第四，在可视化与论证层面，本文已经补充解释方差、Rank IC、滚动稳定性、时间状态、因子热力图、股票潜在结构图、行业聚类交叉图和跨样本边界对比图，使模式发现叙事能够由图表与文字共同支撑。进一步补充的 2015--2026 各年度全 A 长窗口时间状态切换图、因子重要性热力图及年度比较图，已经把此前“长窗口图待补”的占位说明替换为真实证据。当前仍需继续完善的是把这些长窗口结果与更宽输入边界实验联动，而不是继续停留在主线四因子输入上。

针对上述不足，本文给出如下完善计划。其一，在数据层面，继续按照“先明确可用时点，再逐步扩展特征输入”的顺序推进：已完成市场状态代理、外部宏观、财务 PIT、分红和公告事件字典的第一阶段披露时点校验，后续重点转向更长历史、更多宏观品种以及公告正文全文。其二，在评估层面，现有分析已经覆盖前N组合、分组收益、多空、成本后收益、超额收益、回撤、换手、波动、Sharpe 和暴露，后续重点转向字段筛选、长窗口风险归因与交易冲击刻画。其三，在样本层面，保留指数样本池实验作为稳定基线，同时使用全 A 股、行业分层和不同市值分层实验比较不同样本边界下潜在结构、最优秩和稳定性指标的差异。其四，在研究成果展示的相关部分，还可以不断完善长时间统计区间内的各类图表内容，同时也能让绘制完成的图表内容，和正文里面的文字分析内容做到更好的相互呼应。

\chapter{结论与展望}

\section{研究结论}

本文将“股票—因子—时间”三维张量作为核心研究对象，针对股票因子降维与市场模式挖掘的相关研究问题，搭建了规范的数据处理流程、张量建模分析框架，以及结构完整的实验结果输出体系；本文依托HS300、SZ50与ZZ500三类标准样本池的短窗口实验数据，归纳得出了本次研究的主要结论。

首先来看，和股票因子相关的研究内容本身自带多种组合结构，研究人员采用三阶张量的形式整理相关数据，会比普通的二维矩阵形式更加贴合实际的数据特点；这种张量的数据整理方式，既可以留存下个股截面信息，各类因子内容还有时间不断变化之间的关联联系，也能够给之后开展市场规律探寻的相关工作，打下合适的结构条件。

其次，从数据数值复原的层面进行分析，PCA当作二维数据压缩的基础方式，依旧具备较为突出的实用优势；不过在留存原有数据结构，还有解读数据内在规律的层面来看，Tucker分解方式在三组研究样本之内，能够让数据复原效果，个股排序能力以及长期运行稳定程度，达到更为均衡的状态；在Rank IC还有滚动稳定性这类评判指标当中，Tucker对比CP会有着更好的稳定效果，同时这种方式也不会只依靠PCA开展分析，还能够更好助力研究人员完成结构解读和个股排序相关研究工作。

最后，张量分解技术除了能够实现因子数据的降维处理，还可以输出各类隐藏的市场规律信息，包含个股相似结构、因子联动关系以及不同时段的数据状态变动特征；本次实验结果能够看出，质量、价值、动量、波动率等常见因子，在潜在数据空间中具备较强的解释作用，数据整体结构也会在不同时间阶段出现明显的重组变化；这一现象可以说明，张量研究方法能够从数据规律挖掘的角度分析股票因子体系，并不局限于简单的数值压缩功能。

\section{创新点与不足}

本篇论文的研究创新主要体现在三个不同的方面；首先，本文把A股样本实证分析、股票因子降维处理以及市场模式挖掘的相关内容，整合在同一个研究框架中展开分析，能够在一定程度上解决算法分析、数据处理与结论总结相互脱节、各自独立的问题。其次，本文在本科论文的研究基础上，同步对比了CP、Tucker与PCA三种分析方法；对比的内容不局限于数据重构误差这一项指标，还加入了Rank IC与滚动稳定性的对比分析，让最终得出的研究结论更加贴合金融领域的实际研究需求。最后，本文对指数样本、行业样本、市值样本以及长窗口复核实验，搭建了相互对应的对比体系，让论文里的各类图表、数据结果和文字论证内容相互呼应，形成了一套完整统一的实证证据链条。

从整体的研究内容来看，本文没有局限于传统的研究思路，不再只是简单探讨少数因子是否具备有效选股的作用；而是将研究重心前置，重点分析高维因子集合的表征学习与结构压缩相关问题。与此同时，本文也没有将张量分解方法单纯当作压缩数值的数据工具，而是进一步探索了该方法在股票横截面排序、样本边界拓展、行业类别聚类以及时间状态识别等研究场景中的分析解释价值。通过设置这样的研究思路与实验设计，本文能够将多因子模型的选股研究内容，和张量分析的高阶结构研究内容更加自然地结合在一起。

不足之处主要包括：虽然全 A、行业和市值分层实验已经补充到 2015--2026 全部年度的长窗口复核，并形成了相应的跨年度汇总表与比较图，但外部宏观、分红、重大事项和公告文本等补充信息仍以当前实验窗口附近的数据覆盖为主，尚未扩展到更长历史与更完整文本层面；当前实验已经形成前N组合净值、分位数组、多空、成本后净值、基准超额收益和跨样本边界统一汇总，但交易冲击和风险归因仍有继续细化空间；模式发现图表已经从基础指标扩展到股票潜在结构、行业聚类关系和全年长窗口结果，但其中仍有进一步压缩叙述、增强图文配合的空间。

\section{后续工作展望}

后续的研究工作可以从三个维度继续完善推进；第一，在数据研究层面，可进一步扩充目前已接入的外部宏观数据、分红数据、重大事项数据以及公告文本原始数据表，让数据能够覆盖更长的历史区间，收录更完整的公告正文内容，同时细化各类事件的分类标准，全程保障PIT数据的安全性与规范性；第二，在研究方法层面，可继续探索秩自动选择、稀疏约束、非负张量分解以及动态张量分解等相关方法，以此提升模型的解释能力，同时增强模型对不同市场场景的适配效果；第三，在实际应用层面，可将本文现有的结构化实验结果进行延伸拓展，结合长窗口组合复核、回撤风险管控、风险归因分析以及风格轮动监测等相关研究内容，搭建出一套更加完善的量化研究闭环体系。

\chapter*{参考文献}
\addcontentsline{toc}{chapter}{参考文献}
\sloppy
\begin{enumerate}[label={[\arabic*]}]
  \item Hitchcock F L. The Expression of a Tensor or a Polyadic as a Sum of Products[J]. Journal of Mathematical Physics, 1927, 6(1): 164-189.
  \item Tucker L R. Some Mathematical Notes on Three-Mode Factor Analysis[J]. Psychometrika, 1966, 31: 279-311.
  \item Harshman R A. Foundations of the PARAFAC Procedure: Models and Conditions for an Explanatory Multi-modal Factor Analysis[J]. UCLA Working Papers in Phonetics, 1970, 16: 1-84.
  \item Kolda T G, Bader B W. Tensor Decompositions and Applications[J]. SIAM Review, 2009, 51(3): 455-500.
  \item 胡胜龙. 张量分解的唯一性[J]. 运筹学学报(中英文), 2025, 29(03): 34-60.
  \item 夏虹, 张雅倩, 靳晓东, 陈彦萍, 高聪, 王忠民. 基于张量的方法及应用综述[J]. 计算机技术与发展, 2022, 32(06): 1-8.
  \item Rabanser S, Shchur O, Günnemann S. Introduction to Tensor Decompositions and their Applications in Machine Learning[EB/OL]. arXiv:1711.10781, 2017.
  \item 岳欣. 基于张量分解的股票波动分析[D]. 西南财经大学, 2022.
  \item 马少沛, 孙庆慧, 武雅萱, 等. 大数据下张量充分降维方法及其应用研究[J]. 统计研究, 2021, 38(02): 114-134.
  \item 曾亚丽. 基于张量模型的金融时间序列预测研究[D]. 武汉理工大学, 2017.
  \item Brandi G, Gramatica R, Di Matteo T. Unveil Stock Correlation via a New Tensor-based Decomposition Method[EB/OL]. arXiv:1911.06126, 2020.
  \item Spelta A. Financial Market Predictability with Tensor Decomposition and Links Forecast[J]. Applied Network Science, 2017, 2: 7.
  \item 黄宏运, 王梅, 朱家明. 基于多元回归分析的多因子选股模型[J]. 通化师范学院学报, 2016, 37(08): 44-46.
  \item 刘宇轩, 金伟泽, 袁亮. 多因子量化选股模型优化与实证研究——引入金融周期指标的分析[J]. 价格理论与实践, 2022(04): 141-145.
  \item Lettau M. High-Dimensional Factor Models with an Application to Mutual Fund Characteristics[R]. NBER Working Paper No.29833, 2022.
  \item Han Y F, Chen R, Zhang C H. Rank Determination in Tensor Factor Model[EB/OL].\\
  arXiv:2011.07131, 2022.
  \item Wang D, Zheng Y, Lian H, Li G D. High-dimensional Vector Autoregressive Time Series Modeling via Tensor Decomposition[EB/OL]. arXiv:1909.06624, 2020.
  \item 李博, 孟志青, 朱爱花. 时态支持向量机模型在股票操纵模式发现上的研究[J]. 系统科学与数学, 2023, 43(02): 356-378.
  \item 徐常青, 祁力群. 张量CP分解、半正定张量和范德蒙张量[J]. 苏州科技学院学报(自然科学版), 2016, 33(02): 1-7+82.
\end{enumerate}
\fussy

\appendix
\chapter{实验参数与补充说明}

本文正式实验以 HS300、SZ50 与 ZZ500 为样本池。正文首先讨论 2026 年 3 月 2 日至 2026 年 3 月 30 日窗口下的短窗口结果，模型层固定比较 CP、Tucker 与 PCA 三类方法，评价指标包括 MSE、解释方差、Rank IC、IR 与滚动稳定性。实验结果通过统一的结构化评价流程整理，确保论文中的图表、表格和文本结论具有一致证据来源。

\chapter{关键程序代码附录}

为便于说明正文中的关键约束与核心算法，本文选取与主要结论直接相关的程序片段作为附录。附录代码仅保留能够支撑方法解释的关键逻辑，其目的在于说明研究设计如何在实现层面得到落实，而不替代正文中的理论推导与实验分析。

\section{CP 分解的交替最小二乘更新}

下列代码片段对应 CP 分解训练中的核心迭代部分。其基本思路是：在固定另外两个模态因子矩阵后，分别对股票模态、因子模态和时间模态做交替更新，并在每轮迭代后进行列归一化，从而得到可比较的潜在成分权重。

\begin{lstlisting}[language=Python]
for iteration in range(1, max_iter + 1):
    # 固定因子模态与时间模态后更新股票模态载荷。
    gram = (c_mat.T @ c_mat) * (b_mat.T @ b_mat) + identity
    a_mat = (unfold(tensor, 0) @ khatri_rao(c_mat, b_mat)) @ np.linalg.pinv(gram)

    # 固定股票模态与时间模态后更新因子模态载荷。
    gram = (c_mat.T @ c_mat) * (a_mat.T @ a_mat) + identity
    b_mat = (unfold(tensor, 1) @ khatri_rao(c_mat, a_mat)) @ np.linalg.pinv(gram)

    # 固定股票模态与因子模态后更新时间模态载荷。
    gram = (b_mat.T @ b_mat) * (a_mat.T @ a_mat) + identity
    c_mat = (unfold(tensor, 2) @ khatri_rao(b_mat, a_mat)) @ np.linalg.pinv(gram)

    weights = np.ones(rank, dtype=float)
    for factor_matrix in (a_mat, b_mat, c_mat):
        # 列归一化使成分尺度集中到标量权重中，便于比较。
        norms = np.linalg.norm(factor_matrix, axis=0)
        norms[norms == 0] = 1.0
        factor_matrix /= norms
        weights *= norms
\end{lstlisting}

\section{Tucker 分解的多模态交替更新}

下列代码片段体现了 Tucker 分解相对 CP 分解更灵活的地方：不同模态可以设置不同秩，并通过对各模态展开矩阵做奇异向量更新，交替优化股票、因子与时间三个方向的低维表示。核张量 \texttt{core} 用于刻画这些低维表示之间的高阶交互。

\begin{lstlisting}[language=Python]
u_stock = _top_singular_vectors(unfold(tensor, 0), stock_rank)
u_factor = _top_singular_vectors(unfold(tensor, 1), factor_rank)
u_time = _top_singular_vectors(unfold(tensor, 2), time_rank)

for iteration in range(1, max_iter + 1):
    # 在固定因子模态与时间模态后更新股票模态。
    projected = _mode_product(_mode_product(tensor, u_factor.T, 1), u_time.T, 2)
    u_stock = _top_singular_vectors(unfold(projected, 0), stock_rank)

    # 在固定股票模态与时间模态后更新因子模态。
    projected = _mode_product(_mode_product(tensor, u_stock.T, 0), u_time.T, 2)
    u_factor = _top_singular_vectors(unfold(projected, 1), factor_rank)

    # 在固定股票模态与因子模态后更新时间模态。
    projected = _mode_product(_mode_product(tensor, u_stock.T, 0), u_factor.T, 1)
    u_time = _top_singular_vectors(unfold(projected, 2), time_rank)

    # 根据三组模态载荷反推出核张量，保留高阶交互信息。
    core = _mode_product(
        _mode_product(_mode_product(tensor, u_stock.T, 0), u_factor.T, 1),
        u_time.T,
        2,
    )
\end{lstlisting}

\section{PCA 基线的张量展开与重构}

为了与张量方法形成可比的二维压缩基线，本文并未直接在原始表格上做 PCA，而是先把三阶张量沿时间和因子方向展开，再对展开矩阵进行奇异值分解。训练完成后，再把低维结果重构回股票—因子—时间三维张量，以保证评价口径与 CP、Tucker 保持一致。

\begin{lstlisting}[language=Python]
matrix = np.transpose(tensor, (0, 2, 1)).reshape(
    stock_count * time_count,
    factor_count,
)
mean_vector = matrix.mean(axis=0, keepdims=True)
centered = matrix - mean_vector
left, singular_values, right_t = np.linalg.svd(
    centered,
    full_matrices=False,
)

# 以前 rank 个奇异值和奇异向量构造二维低维表示。
scores = left[:, :rank] * singular_values[:rank]
components = right_t[:rank, :].T

# 将二维重构结果恢复到三维张量形状，以保持评价口径一致。
reconstruction_matrix = scores @ components.T + mean_vector
reconstruction = np.transpose(
    reconstruction_matrix.reshape(stock_count, time_count, factor_count),
    (0, 2, 1),
)
\end{lstlisting}

\section{扩展特征的可用时点与时效计算}

下面的代码片段展示了扩展因子面板中“特征年龄”与事件时效衰减的处理方式。其关键点在于：所有时效性特征均从 \texttt{available\_date} 而非 \texttt{pub\_date} 起算，从而保证披露日落在非交易日或收盘后时，模型不会提前获得未来信息。

\begin{lstlisting}[language=Python]
def _snapshot_age_days(snapshot, trade_date: str) -> int:
    if snapshot is None:
        return 0
    # 以最早可交易时点作为年龄起点，保证时点安全。
    effective_date = snapshot.available_date or snapshot.pub_date
    return _days_since(effective_date, trade_date)


def _event_age_decay_score(
    performance_snapshot: PerformanceExpressSnapshot | None,
    forecast_snapshot: ForecastSnapshot | None,
    trade_date: str,
) -> float:
    score = 0.0
    if performance_snapshot is not None:
        # 快报越新，衰减项越大。
        score += 1.0 / (
            1.0 + _snapshot_age_days(performance_snapshot, trade_date)
        )
    if forecast_snapshot is not None:
        # 预告按相同规则计入时间衰减。
        score += 1.0 / (
            1.0 + _snapshot_age_days(forecast_snapshot, trade_date)
        )
    return score
\end{lstlisting}

\section{扩展源表的公告日映射与去重}

下段代码体现了扩展源表构建时的两个关键原则：其一，窗口前事件保持原始公告日，避免历史信息被错误压到首个交易日；其二，公告表按“股票代码—公告类型—公告日—标题—链接”去重，防止同一公告被重复统计到事件特征中。

\begin{lstlisting}[language=Python]
def _source_available_date(pub_date: str, trade_dates: list[str]) -> str:
    # 若公告早于样本窗口的首个交易日，则保留原始日期。
    if trade_dates and pub_date < trade_dates[0]:
        return pub_date
    return _next_trade_date_after(pub_date, trade_dates) or ""


def _notice_identity(
    row: dict[str, object],
) -> tuple[str, str, str, str, str]:
    # 用于识别重复公告，避免同一事件被重复统计到。
    return (
        str(row.get("stock_code", "")),
        str(row.get("notice_type", "")),
        str(row.get("pub_date", "")),
        str(row.get("title_text", "")),
        str(row.get("url", "")),
    )
\end{lstlisting}
